% ============================================================
%  บทที่ 4 — ผลการทดลองเบื้องต้นและระบบต้นแบบ
%  Smart Flow: Digital Twin for KMITL Smart Mobility
%
%  Compiler : XeLaTeX  (Overleaf → Menu → Compiler → XeLaTeX)
%  Encoding : UTF-8
%
%  การใช้งานแบบ standalone  : ไฟล์นี้คอมไพล์ได้เองโดยตรง
%  การใช้งานในเอกสารหลัก   : ตัด preamble ออก แล้วใช้ % ============================================================
%  บทที่ 4 — ผลการทดลองเบื้องต้นและระบบต้นแบบ
%  Smart Flow: Digital Twin for KMITL Smart Mobility
%
%  Compiler : XeLaTeX  (Overleaf → Menu → Compiler → XeLaTeX)
%  Encoding : UTF-8
%
%  การใช้งานแบบ standalone  : ไฟล์นี้คอมไพล์ได้เองโดยตรง
%  การใช้งานในเอกสารหลัก   : ตัด preamble ออก แล้วใช้ % ============================================================
%  บทที่ 4 — ผลการทดลองเบื้องต้นและระบบต้นแบบ
%  Smart Flow: Digital Twin for KMITL Smart Mobility
%
%  Compiler : XeLaTeX  (Overleaf → Menu → Compiler → XeLaTeX)
%  Encoding : UTF-8
%
%  การใช้งานแบบ standalone  : ไฟล์นี้คอมไพล์ได้เองโดยตรง
%  การใช้งานในเอกสารหลัก   : ตัด preamble ออก แล้วใช้ % ============================================================
%  บทที่ 4 — ผลการทดลองเบื้องต้นและระบบต้นแบบ
%  Smart Flow: Digital Twin for KMITL Smart Mobility
%
%  Compiler : XeLaTeX  (Overleaf → Menu → Compiler → XeLaTeX)
%  Encoding : UTF-8
%
%  การใช้งานแบบ standalone  : ไฟล์นี้คอมไพล์ได้เองโดยตรง
%  การใช้งานในเอกสารหลัก   : ตัด preamble ออก แล้วใช้ \input{Chapter4}
% ============================================================

\documentclass[12pt,a4paper]{report}

% ── Thai / Font ─────────────────────────────────────────────
\usepackage{fontspec}
\setmainfont{TH Sarabun New}
\setsansfont{TH Sarabun New}
\setmonofont[Scale=0.82]{Courier New}

% ── Layout ──────────────────────────────────────────────────
\usepackage[
  a4paper,
  left=3.5cm, right=2.5cm,
  top=3.0cm,  bottom=3.0cm
]{geometry}
\usepackage{setspace}
\onehalfspacing
\setlength{\parindent}{1.5em}
\setlength{\parskip}{0.3em}

% ── Math ────────────────────────────────────────────────────
\usepackage{amsmath}
\usepackage{amssymb}

% ── Tables ──────────────────────────────────────────────────
\usepackage{booktabs}
\usepackage{array}
\usepackage{multirow}
\usepackage{makecell}
\usepackage{longtable}
\usepackage{tabularx}
\renewcommand{\arraystretch}{1.5}

% ── Figures ─────────────────────────────────────────────────
\usepackage{graphicx}
\usepackage{float}
\usepackage{caption}
\captionsetup{font=small, labelfont=bf}

% ── Code blocks ─────────────────────────────────────────────
\usepackage{listings}
\usepackage{xcolor}

\definecolor{lstbg}{RGB}{248,248,248}
\definecolor{lstframe}{RGB}{200,200,200}
\definecolor{lstcmt}{RGB}{60,140,60}
\definecolor{lstkw}{RGB}{0,100,160}

\lstset{
  backgroundcolor = \color{lstbg},
  frame           = single,
  rulecolor       = \color{lstframe},
  basicstyle      = \ttfamily\small,
  commentstyle    = \color{lstcmt},
  keywordstyle    = \color{lstkw}\bfseries,
  breaklines      = true,
  numbers         = none,
  captionpos      = b,
  xleftmargin     = 0.8em,
  xrightmargin    = 0.8em,
  aboveskip       = 0.8em,
  belowskip       = 0.6em,
  keepspaces      = true,
  showstringspaces= false,
}

% ── Hyperlinks ───────────────────────────────────────────────
\usepackage[hidelinks]{hyperref}

% ── Custom column types ──────────────────────────────────────
\newcolumntype{C}[1]{>{\centering\arraybackslash}p{#1}}
\newcolumntype{L}[1]{>{\raggedright\arraybackslash}p{#1}}
\newcolumntype{R}[1]{>{\raggedleft\arraybackslash}p{#1}}

% ============================================================
\begin{document}
% ============================================================

% ตั้ง counter ให้บทนี้เป็นบทที่ 4 (หากใช้ \input ในเอกสารหลักให้ลบบรรทัดนี้ออก)
\setcounter{chapter}{3}

\chapter{ผลการทดลองเบื้องต้นและระบบต้นแบบ}
\label{chap:results}

% ============================================================
\section{ผลการทดลองเบื้องต้นและการเลือกองค์ประกอบระบบ}
\label{sec:phase1}
% ============================================================

ในส่วนนี้เป็นการทดลองเบื้องต้นเพื่อคัดเลือกและปรับแต่งค่าพารามิเตอร์ขององค์ประกอบหลัก
ใน Pipeline การประมวลผลวิดีโอ (Video Processing \& AI Pipeline)
เพื่อให้ระบบมีความเหมาะสมและสมดุลที่สุดระหว่างความแม่นยำในการตรวจจับยานพาหนะขนาดเล็ก
กับการทำงานแบบเรียลไทม์

% ────────────────────────────────────────────────────────────
\subsection{การทดสอบขนาดภาพนำเข้าและขีดแบ่งความเชื่อมั่น}
\label{subsec:input-size}
% ────────────────────────────────────────────────────────────

ในการประมวลผลวิดีโอความละเอียด $3840 \times 2160$ พิกเซล จากกล้องมุมสูงบนสะพานลอยถนนฉลองกรุง
พบว่ารถจักรยานยนต์ซึ่งเป็นองค์ประกอบหลักของกระแสจราจรแบบผสม (Mixed Traffic)
ปรากฏเป็นวัตถุขนาดเล็กเมื่ออยู่ไกลจากกล้อง
จึงได้ทำการทดสอบ 2 พารามิเตอร์ที่มีผลต่อการตรวจจับยานพาหนะขนาดเล็กโดยตรง ดังนี้

\paragraph{การทดสอบ Input Image Size}
เปรียบเทียบ imgsz = 640 (ค่ามาตรฐาน YOLOv8) กับ imgsz = 1{,}280
บนวิดีโอตัวอย่าง 525 เฟรม (1920$\times$1080, 25 fps, 21 วินาที)
โดยคงพารามิเตอร์อื่นไว้คงที่ (conf = 0.35, ByteTrack)

\begin{table}[H]
  \centering
  \caption{ผลเปรียบเทียบ Input Image Size ต่อจำนวนยานพาหนะที่นับได้}
  \label{tab:imgsz}
  \small
  \begin{tabular}{L{3.5cm} C{1.8cm} C{2.2cm} C{2.0cm} C{2.0cm} C{2.0cm}}
    \toprule
    \textbf{imgsz} & \textbf{Car} & \textbf{Motorcycle} & \textbf{Truck/Bus} & \textbf{Total} & \textbf{FPS} \\
    \midrule
    640 (ค่ามาตรฐาน) & 5 & \textbf{0} & 0 &  5 & 22.9 \\
    \textbf{1280 (เลือก)} & 3 & \textbf{8} & 0 & 11 & \textbf{18.7} \\
    \bottomrule
  \end{tabular}
\end{table}

\noindent\textbf{วิเคราะห์ผล:}
imgsz = 640 ตรวจจับรถจักรยานยนต์ได้ \textbf{0 คัน} (Recall = 0.000)
เนื่องจากภาพถูกย่อลง 3 เท่าทำให้พิกเซลของรถจักรยานยนต์เล็กเกินขีดจำกัดของโมเดล
imgsz = 1{,}280 ตรวจจับได้ครบ \textbf{8 คัน}
โดยใช้เวลาเพิ่มขึ้นเพียง 5.1 วินาที (จาก 22.9 เป็น 18.7 FPS)
ซึ่งยังอยู่ในระดับที่เหมาะสมสำหรับการทำงานแบบ Real-time

\paragraph{การทดสอบ Confidence Threshold}
ทำการทดสอบค่าขีดแบ่ง 4 ระดับ โดยตรึง imgsz = 1{,}280 ไว้

\begin{table}[H]
  \centering
  \caption{ผลเปรียบเทียบ Confidence Threshold ต่อจำนวน Motorcycle ที่นับได้ (imgsz = 1{,}280)}
  \label{tab:conf}
  \small
  \begin{tabular}{C{3.2cm} C{1.8cm} C{2.5cm} C{2.0cm} C{3.5cm}}
    \toprule
    \textbf{Threshold} & \textbf{Car} & \textbf{Motorcycle} & \textbf{Total} & \textbf{FP (conf $<$ 0.40)} \\
    \midrule
    0.60 (ค่ามาตรฐาน) & 3 & 0            &  3 & 0 \\
    0.50               & 3 & 2            &  5 & 0 \\
    0.40               & 3 & 6            &  9 & 0 \\
    \textbf{0.35 (เลือก)} & 3 & \textbf{8} & \textbf{11} & \textbf{0} \\
    \bottomrule
  \end{tabular}
\end{table}

\noindent\textbf{วิเคราะห์ผล:}
จากการวัดพบว่า Confidence ของรถจักรยานยนต์ ณ เวลาที่ข้ามเส้นนับอยู่ในช่วง
$0.400 \leq \text{conf} \leq 0.549$ (ค่าเฉลี่ย 0.466)
ต่ำกว่ารถยนต์ที่อยู่ในช่วง $0.678 \leq \text{conf} \leq 0.834$ อย่างเป็นระบบ
การใช้ Threshold = 0.60 จึงตัดรถจักรยานยนต์ออกทั้งหมด
ระบบกำหนดค่าขีดแบ่งไว้ที่ \textbf{0.35} ซึ่งคงข้อมูลรถจักรยานยนต์ครบ 8 คัน
โดยไม่มี False Positive เพิ่มเลย เนื่องจากกระบวนการ Vote ของ ByteTrack
กรอง Track ที่ไม่น่าเชื่อถือออกก่อนการนับ

% ────────────────────────────────────────────────────────────
\subsection{การเปรียบเทียบและคัดเลือกอัลกอริทึมติดตามวัตถุ}
\label{subsec:tracking}
% ────────────────────────────────────────────────────────────

ในการทดสอบอัลกอริทึมติดตามวัตถุหลายเป้าหมาย (Multi-Object Tracking)
เพื่อป้องกันการนับยานพาหนะซ้ำ ได้ทำการเปรียบเทียบ 2 อัลกอริทึมหลัก
บนเงื่อนไขเดียวกัน (imgsz = 1{,}280, conf = 0.35, วิดีโอตัวอย่าง 525 เฟรม)

\begin{table}[H]
  \centering
  \caption{ผลเปรียบเทียบ Tracking Algorithm (ByteTrack vs.\ BoT-SORT)}
  \label{tab:tracker}
  \small
  \begin{tabular}{L{6.5cm} C{3.0cm} C{3.0cm}}
    \toprule
    \textbf{ตัวชี้วัด} & \textbf{ByteTrack} & \textbf{BoT-SORT} \\
    \midrule
    Crossing events รวม (คัน)              & 11            & 11       \\
    \quad --- Car                           & 3             & 3        \\
    \quad --- Motorcycle                    & 8             & 8        \\
    ความเร็วประมวลผล (FPS)                  & \textbf{16.4} & 11.5     \\
    Max Motorcycle Track ID (ID Churn)      & 157           & 150      \\
    Anomalies (การนับที่ถูกปฏิเสธ)          & 0             & 0        \\
    Crossing timestamps ตรงกัน             & \checkmark    & \checkmark \\
    \bottomrule
  \end{tabular}
\end{table}

\noindent\textbf{วิเคราะห์ผล:}
ทั้งสองอัลกอริทึมนับยานพาหนะได้จำนวนเท่ากัน (11 คัน)
และตรวจพบยานพาหนะในเวลาเดียวกันทุกรายการ
แสดงให้เห็นว่า Counting Quality ไม่แตกต่างกัน
ByteTrack มีความเร็วสูงกว่า BoT-SORT ถึง \textbf{30\%} (16.4 vs.\ 11.5 FPS)
ในขณะที่ ID Churn ของ Motorcycle ที่ BoT-SORT ต่ำกว่าเพียง 4\%
(Max Track ID 150 vs.\ 157) ซึ่งไม่คุ้มค่ากับ FPS ที่สูญเสียไป
ระบบจึงคัดเลือก \textbf{ByteTrack} เป็นค่าตั้งต้นสำหรับ AI Worker

\medskip
\noindent\textbf{หมายเหตุเกี่ยวกับ ID Churn:}
ค่า Max Motorcycle Track ID สะท้อนจำนวน Track ทั้งหมดที่ถูกสร้างขึ้นตลอดคลิป
(ByteTrack สร้าง 157 Track สำหรับรถจักรยานยนต์ที่ข้ามเส้น 8 คัน = อัตราส่วน 19.6:1)
ค่าที่สูงนี้เป็นปกติสำหรับจราจรแบบผสม เนื่องจาก Tracker สร้าง Track ชั่วคราว
สำหรับยานพาหนะที่เข้ามาในโซนแต่ยังไม่ถึงเส้นนับ
ByteTrack จัดการ Track เหล่านี้ด้วย Two-stage Data Association
ซึ่งฟื้นฟู Track ที่หายไปชั่วคราวได้ โดยไม่เพิ่มจำนวน Crossing Event ที่ผิดพลาด

% ────────────────────────────────────────────────────────────
\subsection{การกำหนดขนาดหน้าต่างเวลาประเมินสภาพจราจร}
\label{subsec:time-window}
% ────────────────────────────────────────────────────────────

ระบบทำการทดสอบความยาวหน้าต่างเวลา ($T$) ในการคำนวณอัตราการไหล ($V$)
และค่าความหนาแน่น ($O$)
พบว่าหน้าต่างเวลาที่ \textbf{10 วินาที} ให้ค่าเฉลี่ยสถิติการข้ามเส้น
ในระดับหลักหน่วยถึงหลักสิบ
ซึ่งไม่ผันผวนรวดเร็วเกินไปจนระบบแกว่ง (เมื่อเทียบกับ 5 วินาที)
และไม่ล่าช้าเกินไปในการตอบสนองต่อสภาวะรถติด (เมื่อเทียบกับ 30 วินาที)

% ============================================================
\section{ผลการทดสอบประสิทธิภาพของระบบและการจำแนกสถานะการจราจร}
\label{sec:phase2}
% ============================================================

ผลการทดลองทั้ง 4 ชุดตามที่ออกแบบไว้ในแผนงาน
เพื่อประเมินความถูกต้องของการนับ การจำแนกสภาวะ
ผลของการแยกอัตราการไหล และความหน่วงของข้อมูล สรุปได้ดังนี้

% ────────────────────────────────────────────────────────────
\subsection{ผลการทดลองที่ 1: ความแม่นยำของการตรวจจับและนับยานพาหนะ}
\label{subsec:exp1}
% ────────────────────────────────────────────────────────────

ทำการเปรียบเทียบยอดนับจากระบบ AI Worker
กับค่าอ้างอิงจากการนับด้วยตาเปล่า (Ground Truth) ในชุดข้อมูลทดสอบ
โดยแยกประเมินตามประเภทยานพาหนะ
เพื่อทดสอบสมมติฐานด้านความยากในการตรวจจับรถจักรยานยนต์

\begin{table}[H]
  \centering
  \caption{ผลการทดสอบความแม่นยำของการตรวจจับและนับยานพาหนะแยกตามประเภท}
  \label{tab:exp1}
  \small
  \setlength{\tabcolsep}{5pt}
  \begin{tabular}{L{3.0cm} C{1.5cm} C{1.6cm} C{0.8cm} C{0.8cm} C{1.4cm} C{1.3cm} C{1.5cm} C{1.4cm}}
    \toprule
    \textbf{ประเภทยานพาหนะ} &
    \makecell{\textbf{Ground}\\\textbf{Truth}\\\textbf{(คัน)}} &
    \makecell{\textbf{System}\\\textbf{Count}\\\textbf{(TP)}} &
    \textbf{FP} &
    \textbf{FN} &
    \textbf{Precision} &
    \textbf{Recall} &
    \textbf{F1-Score} &
    \makecell{\textbf{MAE}\\\textbf{(คัน)}} \\
    \midrule
    รถยนต์ (Car)                       & 18          & 17          & 1 & 1 & 0.944 & 0.944 & 0.944 & 0.056 \\
    รถบรรทุก/รถบัส (Truck/Bus)         & 2           & 2           & 0 & 0 & 1.000 & 1.000 & 1.000 & 0.000 \\
    รถจักรยานยนต์ (Motorcycle)         & 8           & 7           & 1 & 1 & 0.875 & 0.875 & 0.875 & 0.125 \\
    \midrule
    \textbf{ภาพรวมทุกประเภท (Overall)} & \textbf{28} & \textbf{26} & \textbf{2} & \textbf{2}
        & \textbf{0.929} & \textbf{0.929} & \textbf{0.929} & \textbf{0.061} \\
    \bottomrule
  \end{tabular}
\end{table}

\noindent\textbf{วิเคราะห์ผล:}
ค่า Recall ของรถยนต์อยู่ที่ 0.944 (ผ่านเกณฑ์ $\geq 0.85$)
ค่า Recall ของรถจักรยานยนต์อยู่ที่ 0.875 (ผ่านเกณฑ์ $\geq 0.70$)
และค่า Precision ภาพรวมเท่ากับ 0.929 (ผ่านเกณฑ์ $\geq 0.85$)
เหตุผลที่รถจักรยานยนต์มี FN เกิดจากกรณีที่โมเดลจำแนกผิดเป็นชนิดจักรยาน (Bicycle)
ซึ่งถูกตั้งใจตัดออกตั้งแต่ชั้น AI Worker

% ────────────────────────────────────────────────────────────
\subsection{ผลการทดลองที่ 2: ความแม่นยำของการจำแนกสภาพการจราจร 4 สถานะ}
\label{subsec:exp2}
% ────────────────────────────────────────────────────────────

ประเมินการตัดสินสภาวะการจราจรในแต่ละหน้าต่างเวลา 10 วินาที
เทียบกับการกำกับสถานะจากวิดีโอจริง

\begin{table}[H]
  \centering
  \caption{Confusion Matrix การจำแนกสภาพการจราจร 4 สถานะ}
  \label{tab:confusion}
  \small
  \setlength{\tabcolsep}{4pt}
  \begin{tabular}{L{3.8cm} C{1.4cm} C{1.8cm} C{2.0cm} C{1.7cm} C{1.4cm} C{1.3cm}}
    \toprule
    \makecell[l]{\textbf{สภาพจริง} \\ $\backslash$\,\textbf{ผลประเมิน}} &
    \textbf{Normal} &
    \makecell{\textbf{High}\\\textbf{Density}} &
    \makecell{\textbf{Slow}\\\textbf{Moving}} &
    \textbf{Standstill} &
    \textbf{Precision} &
    \textbf{Recall} \\
    \midrule
    Normal ($O < 3.0$)
        & \textbf{12} & 0          & 0          & 0          & 1.000 & 1.000 \\
    High Density ($O \geq 3.0,\, V_c \geq 12.0$)
        & 1           & \textbf{4} & 0          & 0          & 1.000 & 0.800 \\
    Slow Moving ($O \geq 3.0,\, 2.0 \leq V_c < 12.0$)
        & 0           & 0          & \textbf{3} & 0          & 1.000 & 1.000 \\
    Standstill ($O \geq 3.0,\, V_c < 2.0$)
        & 0           & 0          & 0          & \textbf{2} & 1.000 & 1.000 \\
    \bottomrule
  \end{tabular}
\end{table}

\noindent\textbf{วิเคราะห์ผล:}
ความถูกต้องโดยรวม (Overall Accuracy) เท่ากับ \textbf{0.954 (21/22 windows)}
ซึ่งสูงกว่าเกณฑ์ที่กำหนดไว้ ($\geq 0.70$)
และที่สำคัญที่สุดคือ
\textbf{ไม่พบกรณีการตัดสินสภาวะหยุดนิ่ง (Standstill) เป็นการจราจรปกติ (Normal)}
ซึ่งสอดคล้องกับเกณฑ์การยอมรับขั้นเด็ดขาด (Must not occur)

% ────────────────────────────────────────────────────────────
\subsection{ผลการทดลองที่ 3: ผลของการแยกอัตราการไหลของรถจักรยานยนต์}
\label{subsec:exp3}
% ────────────────────────────────────────────────────────────

ทดสอบข้อเสนอหลักของโครงงาน
โดยเปรียบเทียบคำตัดสินสภาพการจราจรระหว่าง
\textbf{วิธีที่แยกอัตราการไหลของรถจักรยานยนต์ ($V_c$)}
กับ
\textbf{วิธีดั้งเดิมที่นับรวมยานพาหนะทุกประเภท ($V_{total}$)}

ข้อมูลทดสอบที่ใช้สาธิต (ช่วงวินาทีที่ 10--20):

\begin{lstlisting}[caption={ค่าพารามิเตอร์สภาพจราจรช่วงวินาทีที่ 10--20 ในการทดสอบ Experiment 3}]
ความหนาแน่นเฉลี่ย (O)          : 16.85 คัน/เฟรม  (รถยนต์จอดนิ่งเต็มโซน)
อัตราการไหลรถยนต์ (Vc)         : 0.0 คัน/นาที    (ไม่มีรถยนต์ผ่านเส้นนับ)
อัตราการไหลรถจักรยานยนต์ (Vmc) : 24.0 คัน/นาที  (รถจักรยานยนต์แทรกผ่านได้ 4 คัน)
\end{lstlisting}

\noindent\textbf{ผลการตัดสิน:}

\begin{itemize}
  \item \textbf{วิธีดั้งเดิม (นับรวม $V_{total} = 24.0$ คัน/นาที):}
        ระบบตัดสินเป็น \textbf{High Density}
        (จราจรหนาแน่นแต่ยังไหลได้ดี)
        ซึ่ง\textbf{ผิดจากความเป็นจริง} เพราะรถยนต์หยุดนิ่งสนิท

  \item \textbf{วิธีของโครงงาน (แยกเฉพาะ $V_c = 0.0$ คัน/นาที):}
        ระบบตัดสินเป็น \textbf{Standstill}
        (จราจรติดขัดหยุดนิ่ง)
        ซึ่ง\textbf{ถูกต้องตรงตามสภาพจริง}
\end{itemize}

\noindent\textbf{สรุป:}
มีหน้าต่างเวลาที่คำตัดสินเปลี่ยนไปอย่างมีนัยสำคัญตามเกณฑ์การยอมรับ
($\geq 1$ หน้าต่าง)
พิสูจน์ให้เห็นว่าการแยก $V_{mc}$ ออกจากเกณฑ์ตัดสิน
ช่วยป้องกันไม่ให้รถจักรยานยนต์กลบสัญญาณสภาวะรถติดรุนแรงได้จริง

% ────────────────────────────────────────────────────────────
\subsection{ผลการทดลองที่ 4: ความหน่วงของการเชื่อมโยงข้อมูล}
\label{subsec:exp4}
% ────────────────────────────────────────────────────────────

ทำการแทรกจุดประทับเวลา (Timestamping) 4 จุด
เพื่อวัดค่าความหน่วง End-to-End จาก AI Worker ถึงการแสดงผลบน Unity

\begin{table}[H]
  \centering
  \caption{ผลการวัดความหน่วงของระบบแบ่งตามช่วงการส่งผ่านข้อมูล}
  \label{tab:latency}
  \begin{tabular}{L{8.2cm} C{3.0cm} C{3.0cm}}
    \toprule
    \textbf{ช่วงการส่งผ่านข้อมูล (Processing Pipeline Segment)} &
    \textbf{Median Latency (ms)} &
    \textbf{P95 Latency (ms)} \\
    \midrule
    1.\ AI Worker Processing \& Event Generation              & 22.4          & 38.1 \\
    2.\ HTTP POST Ingest to Backend Validation                & 4.8           & 11.2 \\
    3.\ Backend WebSocket Broadcast (\texttt{spawn\_vehicle}) & 2.1           & 5.5  \\
    4.\ Unity WebSocket Receive to Object Pool Spawn          & 18.5          & 45.0 \\
    \midrule
    \textbf{รวมความหน่วงข้อมูลรายคัน (End-to-End Latency)}   & \textbf{47.8} & \textbf{99.8} \\
    \bottomrule
  \end{tabular}
\end{table}

\noindent\textbf{วิเคราะห์ผล:}
ความหน่วงรวมเฉลี่ย (Median Latency) ของเหตุการณ์รายคันเท่ากับ
\textbf{47.8 มิลลิวินาที (0.0478 วินาที)}
ซึ่งดีกว่าเกณฑ์ที่กำหนดไว้ไม่เกิน 1 วินาทีอย่างมาก
สำหรับคิวการสร้างใน Unity ไม่เกิดปัญหาสะสมคิวเนื่องจากอัตราการไหลของรถในวิดีโอ
ไม่เกินขีดจำกัดการดึงวัตถุจาก Object Pool (2 คัน/วินาที)

% ============================================================
\section{รายละเอียดและการทำงานของระบบต้นแบบ}
\label{sec:prototype}
% ============================================================

สถาปัตยกรรมระบบต้นแบบแฝดดิจิทัล (Digital Shadow level)
ประกอบด้วย 3 โมดูลหลักที่ทำงานร่วมกันอย่างเป็นระบบ
ดังแผนภาพในรูปที่~\ref{fig:pipeline}

\begin{figure}[H]
\begin{lstlisting}[language=,frame=single,basicstyle=\ttfamily\scriptsize,
                   caption={สถาปัตยกรรม End-to-End Pipeline ของระบบต้นแบบ Smart Flow},
                   label={fig:pipeline}]
+----------------------------------------------------------------------+
|                           PHYSICAL LAYER                             |
|           กล้อง CCTV สะพานลอย ถนนฉลองกรุง (3840x2160, 30fps)        |
+----------------------------------------------------------------------+
                                  | (Video Stream)
                                  v
+----------------------------------------------------------------------+
|               AI WORKER  (Python / YOLOv8s / ByteTrack)             |
|  - ตรวจจับ/ติดตามยานพาหนะ  (ImgSize=1280, Conf=0.35)                |
|  - คำนวณ Line Crossing & Density Polygon                             |
|  - สรุป Traffic State ทุก 10 วินาที (แยก Vc และ Vmc)                |
+----------------------------------------------------------------------+
         | HTTP POST  (/api/ingest  per vehicle, /api/traffic-state)
                                  v
+----------------------------------------------------------------------+
|            BACKEND SERVER  (Node.js / Express / Socket.IO)          |
|  - Data Contract Enforcement & HTTP 400 Validation                  |
|  - In-Memory Counter Accumulation                                   |
|  - WebSocket Broadcasting  ('spawn_vehicle', 'traffic_state')       |
+----------------------------------------------------------------------+
                       | WebSocket  (Socket.IO JSON)
                                  v
+----------------------------------------------------------------------+
|                  VIRTUAL LAYER  (Unity 3D / C# / WebGL)             |
|  - Event Listener & Queue Manager (FIFO)                            |
|  - Object Pooling Vehicle Spawning & Waypoint Navigation            |
|  - Real-time 3D Scene Environment Sync                              |
+----------------------------------------------------------------------+
\end{lstlisting}
\end{figure}

% ────────────────────────────────────────────────────────────
\subsection{ส่วนประสานผู้ใช้และการแสดงผลสามมิติ (Unity 3D Virtual Layer)}
\label{subsec:unity}
% ────────────────────────────────────────────────────────────

ส่วนแสดงผล 3D พัฒนาบน Unity เอนจินและปรับแต่งสำหรับ WebGL/PC
ประกอบด้วยองค์ประกอบหลัก 3 ส่วน

\begin{enumerate}
  \item \textbf{3D Virtual City View:}
        จำลองโครงสร้างทางกายภาพของถนนฉลองกรุงและสิ่งปลูกสร้างโดยรอบ
        สปอว์นโมเดล 3D ของรถยนต์ Bus Truck และ Motorcycle
        เคลื่อนที่ตามจุดนำทาง (Waypoints) บนเลนที่ถูกต้องแบบเรียลไทม์

  \item \textbf{Object Pooling System:}
        สร้าง Prefab ยานพาหนะสำรองไว้ในหน่วยความจำล่วงหน้า
        เมื่อได้รับอีเวนต์ \texttt{spawn\_vehicle}
        จะดึงวัตถุมาเปิดใช้งาน (Enable)
        และคืนค่าเข้าคิว (Disable) เมื่อวิ่งสุดเส้นทาง
        ช่วยลดอาการกระตุกของ Frame Rate บน Unity

  \item \textbf{Developer Monitoring Dashboard:}
        หน้าเว็บสำหรับผู้พัฒนาที่รับข้อมูลผ่าน WebSocket ช่องทาง
        \texttt{spawn} และ \texttt{traffic\_state}
        แสดงสถานะตัวนับสะสม (Counter Statistics)
        และผลตัดสินสภาพการจราจรในรูปแบบตัวเลขและกราฟแบบ Real-time
\end{enumerate}

% ────────────────────────────────────────────────────────────
\subsection{สภาพการณ์การใช้งานระบบ (Use Case Scenarios)}
\label{subsec:scenarios}
% ────────────────────────────────────────────────────────────

ระบบรองรับการทำงานในสถานการณ์จริง 4 รูปแบบหลัก

\begin{itemize}
  \item \textbf{Scenario 1 --- Normal Traffic Flow Sync:}
        ยานพาหนะข้ามเส้นนับ
        AI Worker ส่งเหตุการณ์รายคันผ่าน HTTP POST
        Backend ตรวจสอบผ่าน Data Contract แล้วส่งต่อผ่าน WebSocket
        ให้ Unity สปอว์นโมเดลรถวิ่งด้วยความเร็วปกติทันที

  \item \textbf{Scenario 2 --- Standstill Detection during Mixed Traffic:}
        เกิดสภาวะรถยนต์ติดสะสมหยุดนิ่ง แต่รถจักรยานยนต์ยังแทรกตัวผ่านได้
        AI Worker สรุปค่า $O = 16.85$ และ $V_c = 0.0$
        ส่งคำตัดสิน \texttt{standstill} ไปยัง Backend
        และกระจายไปปรับสภาวะฉาก 3D

  \item \textbf{Scenario 3 --- Invalid Data Handling (Data Contract Gate):}
        หากโมดูลต้นทางส่งข้อมูลที่มีโครงสร้างผิดพลาด
        เช่น \texttt{confidence} ไม่อยู่ในช่วง 0.0--1.0 หรือเวลาติดลบ
        Backend จะปฏิเสธคำขอด้วยรหัส \texttt{HTTP 400 Bad Request} ทันที
        ป้องกันไม่ให้ข้อมูลขยะหลุดไปสร้างความผิดพลาดในส่วนแสดงผล 3D

  \item \textbf{Scenario 4 --- Track Loss \& Auto-Recovery:}
        หากยานพาหนะถูกบดบังชั่วขณะ
        อัลกอริทึม ByteTrack จะรักษาสถานะ Tracklet ไว้ในคิว Lost Track
        เป็นเวลา 30 เฟรม ก่อนจะจับคู่คืนรหัส ID เดิมเมื่อวัตถุปรากฏขึ้นอีกครั้ง
        ป้องกันการสร้างยานพาหนะซ้ำซ้อนในโลกเสมือน
\end{itemize}

% ============================================================
\section{สรุปการประเมินผลเทียบกับเกณฑ์การยอมรับ}
\label{sec:summary}
% ============================================================

จากการทดสอบระบบตามตัวชี้วัดทั้งหมดที่กำหนดไว้ในบทที่ 3
สามารถสรุปผลการประเมินเปรียบเทียบได้ดังตารางที่~\ref{tab:acceptance}

\begin{table}[H]
  \centering
  \caption{สรุปผลการประเมินเทียบกับเกณฑ์การยอมรับทั้งหมด}
  \label{tab:acceptance}
  \small
  \setlength{\tabcolsep}{5pt}
  \begin{tabular}{L{4.5cm} C{3.2cm} C{3.0cm} C{2.5cm}}
    \toprule
    \textbf{การทดลอง / ตัวชี้วัด} &
    \makecell{\textbf{เกณฑ์การยอมรับ}\\\textbf{(Target)}} &
    \makecell{\textbf{ผลที่ได้}\\\textbf{(Actual Result)}} &
    \textbf{การประเมิน} \\
    \midrule
    Exp 1: Recall ของรถยนต์ (Car)
        & $\geq 0.85$
        & \textbf{0.944}
        & \textbf{ผ่านเกณฑ์} \\
    Exp 1: Recall ของรถจักรยานยนต์ (Motorcycle)
        & $\geq 0.70$
        & \textbf{0.875}
        & \textbf{ผ่านเกณฑ์} \\
    Exp 1: Precision รวมทุกประเภท
        & $\geq 0.85$
        & \textbf{0.929}
        & \textbf{ผ่านเกณฑ์} \\
    Exp 2: ความถูกต้องรวมการจำแนกสถานะ
        & $\geq 0.70$
        & \textbf{0.954}
        & \textbf{ผ่านเกณฑ์} \\
    Exp 2: ตัดสิน Standstill เป็น Normal
        & ต้องไม่เกิดขึ้นเลย
        & \textbf{0 กรณี}
        & \textbf{ผ่านเกณฑ์} \\
    Exp 3: หน้าต่างที่คำตัดสินเปลี่ยน
        & $\geq 1$ หน้าต่าง
        & \textbf{1 หน้าต่าง}
        & \textbf{ผ่านเกณฑ์} \\
    Exp 4: ความหน่วงข้อมูลรายคัน (Median)
        & $\leq 1.0$ วินาที
        & \textbf{0.0478 วินาที}
        & \textbf{ผ่านเกณฑ์} \\
    \bottomrule
  \end{tabular}
\end{table}

ผลการทดลองในบทนี้ยืนยันว่าระบบต้นแบบแฝดดิจิทัลที่พัฒนาขึ้น
สามารถประมวลผลวิดีโอ แยกประเภทยานพาหนะ
ตัดสินสภาพการจราจรแบบผสมในประเทศไทยได้อย่างแม่นยำ
และส่งต่อข้อมูลไปจำลองภาพ 3D แบบเรียลไทม์ได้ตรงตามเป้าหมายของโครงงานทุกประการ

% ============================================================
\end{document}
% ============================================================

% ============================================================

\documentclass[12pt,a4paper]{report}

% ── Thai / Font ─────────────────────────────────────────────
\usepackage{fontspec}
\setmainfont{TH Sarabun New}
\setsansfont{TH Sarabun New}
\setmonofont[Scale=0.82]{Courier New}

% ── Layout ──────────────────────────────────────────────────
\usepackage[
  a4paper,
  left=3.5cm, right=2.5cm,
  top=3.0cm,  bottom=3.0cm
]{geometry}
\usepackage{setspace}
\onehalfspacing
\setlength{\parindent}{1.5em}
\setlength{\parskip}{0.3em}

% ── Math ────────────────────────────────────────────────────
\usepackage{amsmath}
\usepackage{amssymb}

% ── Tables ──────────────────────────────────────────────────
\usepackage{booktabs}
\usepackage{array}
\usepackage{multirow}
\usepackage{makecell}
\usepackage{longtable}
\usepackage{tabularx}
\renewcommand{\arraystretch}{1.5}

% ── Figures ─────────────────────────────────────────────────
\usepackage{graphicx}
\usepackage{float}
\usepackage{caption}
\captionsetup{font=small, labelfont=bf}

% ── Code blocks ─────────────────────────────────────────────
\usepackage{listings}
\usepackage{xcolor}

\definecolor{lstbg}{RGB}{248,248,248}
\definecolor{lstframe}{RGB}{200,200,200}
\definecolor{lstcmt}{RGB}{60,140,60}
\definecolor{lstkw}{RGB}{0,100,160}

\lstset{
  backgroundcolor = \color{lstbg},
  frame           = single,
  rulecolor       = \color{lstframe},
  basicstyle      = \ttfamily\small,
  commentstyle    = \color{lstcmt},
  keywordstyle    = \color{lstkw}\bfseries,
  breaklines      = true,
  numbers         = none,
  captionpos      = b,
  xleftmargin     = 0.8em,
  xrightmargin    = 0.8em,
  aboveskip       = 0.8em,
  belowskip       = 0.6em,
  keepspaces      = true,
  showstringspaces= false,
}

% ── Hyperlinks ───────────────────────────────────────────────
\usepackage[hidelinks]{hyperref}

% ── Custom column types ──────────────────────────────────────
\newcolumntype{C}[1]{>{\centering\arraybackslash}p{#1}}
\newcolumntype{L}[1]{>{\raggedright\arraybackslash}p{#1}}
\newcolumntype{R}[1]{>{\raggedleft\arraybackslash}p{#1}}

% ============================================================
\begin{document}
% ============================================================

% ตั้ง counter ให้บทนี้เป็นบทที่ 4 (หากใช้ \input ในเอกสารหลักให้ลบบรรทัดนี้ออก)
\setcounter{chapter}{3}

\chapter{ผลการทดลองเบื้องต้นและระบบต้นแบบ}
\label{chap:results}

% ============================================================
\section{ผลการทดลองเบื้องต้นและการเลือกองค์ประกอบระบบ}
\label{sec:phase1}
% ============================================================

ในส่วนนี้เป็นการทดลองเบื้องต้นเพื่อคัดเลือกและปรับแต่งค่าพารามิเตอร์ขององค์ประกอบหลัก
ใน Pipeline การประมวลผลวิดีโอ (Video Processing \& AI Pipeline)
เพื่อให้ระบบมีความเหมาะสมและสมดุลที่สุดระหว่างความแม่นยำในการตรวจจับยานพาหนะขนาดเล็ก
กับการทำงานแบบเรียลไทม์

% ────────────────────────────────────────────────────────────
\subsection{การทดสอบขนาดภาพนำเข้าและขีดแบ่งความเชื่อมั่น}
\label{subsec:input-size}
% ────────────────────────────────────────────────────────────

ในการประมวลผลวิดีโอความละเอียด $3840 \times 2160$ พิกเซล จากกล้องมุมสูงบนสะพานลอยถนนฉลองกรุง
พบว่ารถจักรยานยนต์ซึ่งเป็นองค์ประกอบหลักของกระแสจราจรแบบผสม (Mixed Traffic)
ปรากฏเป็นวัตถุขนาดเล็กเมื่ออยู่ไกลจากกล้อง
จึงได้ทำการทดสอบ 2 พารามิเตอร์ที่มีผลต่อการตรวจจับยานพาหนะขนาดเล็กโดยตรง ดังนี้

\paragraph{การทดสอบ Input Image Size}
เปรียบเทียบ imgsz = 640 (ค่ามาตรฐาน YOLOv8) กับ imgsz = 1{,}280
บนวิดีโอตัวอย่าง 525 เฟรม (1920$\times$1080, 25 fps, 21 วินาที)
โดยคงพารามิเตอร์อื่นไว้คงที่ (conf = 0.35, ByteTrack)

\begin{table}[H]
  \centering
  \caption{ผลเปรียบเทียบ Input Image Size ต่อจำนวนยานพาหนะที่นับได้}
  \label{tab:imgsz}
  \small
  \begin{tabular}{L{3.5cm} C{1.8cm} C{2.2cm} C{2.0cm} C{2.0cm} C{2.0cm}}
    \toprule
    \textbf{imgsz} & \textbf{Car} & \textbf{Motorcycle} & \textbf{Truck/Bus} & \textbf{Total} & \textbf{FPS} \\
    \midrule
    640 (ค่ามาตรฐาน) & 5 & \textbf{0} & 0 &  5 & 22.9 \\
    \textbf{1280 (เลือก)} & 3 & \textbf{8} & 0 & 11 & \textbf{18.7} \\
    \bottomrule
  \end{tabular}
\end{table}

\noindent\textbf{วิเคราะห์ผล:}
imgsz = 640 ตรวจจับรถจักรยานยนต์ได้ \textbf{0 คัน} (Recall = 0.000)
เนื่องจากภาพถูกย่อลง 3 เท่าทำให้พิกเซลของรถจักรยานยนต์เล็กเกินขีดจำกัดของโมเดล
imgsz = 1{,}280 ตรวจจับได้ครบ \textbf{8 คัน}
โดยใช้เวลาเพิ่มขึ้นเพียง 5.1 วินาที (จาก 22.9 เป็น 18.7 FPS)
ซึ่งยังอยู่ในระดับที่เหมาะสมสำหรับการทำงานแบบ Real-time

\paragraph{การทดสอบ Confidence Threshold}
ทำการทดสอบค่าขีดแบ่ง 4 ระดับ โดยตรึง imgsz = 1{,}280 ไว้

\begin{table}[H]
  \centering
  \caption{ผลเปรียบเทียบ Confidence Threshold ต่อจำนวน Motorcycle ที่นับได้ (imgsz = 1{,}280)}
  \label{tab:conf}
  \small
  \begin{tabular}{C{3.2cm} C{1.8cm} C{2.5cm} C{2.0cm} C{3.5cm}}
    \toprule
    \textbf{Threshold} & \textbf{Car} & \textbf{Motorcycle} & \textbf{Total} & \textbf{FP (conf $<$ 0.40)} \\
    \midrule
    0.60 (ค่ามาตรฐาน) & 3 & 0            &  3 & 0 \\
    0.50               & 3 & 2            &  5 & 0 \\
    0.40               & 3 & 6            &  9 & 0 \\
    \textbf{0.35 (เลือก)} & 3 & \textbf{8} & \textbf{11} & \textbf{0} \\
    \bottomrule
  \end{tabular}
\end{table}

\noindent\textbf{วิเคราะห์ผล:}
จากการวัดพบว่า Confidence ของรถจักรยานยนต์ ณ เวลาที่ข้ามเส้นนับอยู่ในช่วง
$0.400 \leq \text{conf} \leq 0.549$ (ค่าเฉลี่ย 0.466)
ต่ำกว่ารถยนต์ที่อยู่ในช่วง $0.678 \leq \text{conf} \leq 0.834$ อย่างเป็นระบบ
การใช้ Threshold = 0.60 จึงตัดรถจักรยานยนต์ออกทั้งหมด
ระบบกำหนดค่าขีดแบ่งไว้ที่ \textbf{0.35} ซึ่งคงข้อมูลรถจักรยานยนต์ครบ 8 คัน
โดยไม่มี False Positive เพิ่มเลย เนื่องจากกระบวนการ Vote ของ ByteTrack
กรอง Track ที่ไม่น่าเชื่อถือออกก่อนการนับ

% ────────────────────────────────────────────────────────────
\subsection{การเปรียบเทียบและคัดเลือกอัลกอริทึมติดตามวัตถุ}
\label{subsec:tracking}
% ────────────────────────────────────────────────────────────

ในการทดสอบอัลกอริทึมติดตามวัตถุหลายเป้าหมาย (Multi-Object Tracking)
เพื่อป้องกันการนับยานพาหนะซ้ำ ได้ทำการเปรียบเทียบ 2 อัลกอริทึมหลัก
บนเงื่อนไขเดียวกัน (imgsz = 1{,}280, conf = 0.35, วิดีโอตัวอย่าง 525 เฟรม)

\begin{table}[H]
  \centering
  \caption{ผลเปรียบเทียบ Tracking Algorithm (ByteTrack vs.\ BoT-SORT)}
  \label{tab:tracker}
  \small
  \begin{tabular}{L{6.5cm} C{3.0cm} C{3.0cm}}
    \toprule
    \textbf{ตัวชี้วัด} & \textbf{ByteTrack} & \textbf{BoT-SORT} \\
    \midrule
    Crossing events รวม (คัน)              & 11            & 11       \\
    \quad --- Car                           & 3             & 3        \\
    \quad --- Motorcycle                    & 8             & 8        \\
    ความเร็วประมวลผล (FPS)                  & \textbf{16.4} & 11.5     \\
    Max Motorcycle Track ID (ID Churn)      & 157           & 150      \\
    Anomalies (การนับที่ถูกปฏิเสธ)          & 0             & 0        \\
    Crossing timestamps ตรงกัน             & \checkmark    & \checkmark \\
    \bottomrule
  \end{tabular}
\end{table}

\noindent\textbf{วิเคราะห์ผล:}
ทั้งสองอัลกอริทึมนับยานพาหนะได้จำนวนเท่ากัน (11 คัน)
และตรวจพบยานพาหนะในเวลาเดียวกันทุกรายการ
แสดงให้เห็นว่า Counting Quality ไม่แตกต่างกัน
ByteTrack มีความเร็วสูงกว่า BoT-SORT ถึง \textbf{30\%} (16.4 vs.\ 11.5 FPS)
ในขณะที่ ID Churn ของ Motorcycle ที่ BoT-SORT ต่ำกว่าเพียง 4\%
(Max Track ID 150 vs.\ 157) ซึ่งไม่คุ้มค่ากับ FPS ที่สูญเสียไป
ระบบจึงคัดเลือก \textbf{ByteTrack} เป็นค่าตั้งต้นสำหรับ AI Worker

\medskip
\noindent\textbf{หมายเหตุเกี่ยวกับ ID Churn:}
ค่า Max Motorcycle Track ID สะท้อนจำนวน Track ทั้งหมดที่ถูกสร้างขึ้นตลอดคลิป
(ByteTrack สร้าง 157 Track สำหรับรถจักรยานยนต์ที่ข้ามเส้น 8 คัน = อัตราส่วน 19.6:1)
ค่าที่สูงนี้เป็นปกติสำหรับจราจรแบบผสม เนื่องจาก Tracker สร้าง Track ชั่วคราว
สำหรับยานพาหนะที่เข้ามาในโซนแต่ยังไม่ถึงเส้นนับ
ByteTrack จัดการ Track เหล่านี้ด้วย Two-stage Data Association
ซึ่งฟื้นฟู Track ที่หายไปชั่วคราวได้ โดยไม่เพิ่มจำนวน Crossing Event ที่ผิดพลาด

% ────────────────────────────────────────────────────────────
\subsection{การกำหนดขนาดหน้าต่างเวลาประเมินสภาพจราจร}
\label{subsec:time-window}
% ────────────────────────────────────────────────────────────

ระบบทำการทดสอบความยาวหน้าต่างเวลา ($T$) ในการคำนวณอัตราการไหล ($V$)
และค่าความหนาแน่น ($O$)
พบว่าหน้าต่างเวลาที่ \textbf{10 วินาที} ให้ค่าเฉลี่ยสถิติการข้ามเส้น
ในระดับหลักหน่วยถึงหลักสิบ
ซึ่งไม่ผันผวนรวดเร็วเกินไปจนระบบแกว่ง (เมื่อเทียบกับ 5 วินาที)
และไม่ล่าช้าเกินไปในการตอบสนองต่อสภาวะรถติด (เมื่อเทียบกับ 30 วินาที)

% ============================================================
\section{ผลการทดสอบประสิทธิภาพของระบบและการจำแนกสถานะการจราจร}
\label{sec:phase2}
% ============================================================

ผลการทดลองทั้ง 4 ชุดตามที่ออกแบบไว้ในแผนงาน
เพื่อประเมินความถูกต้องของการนับ การจำแนกสภาวะ
ผลของการแยกอัตราการไหล และความหน่วงของข้อมูล สรุปได้ดังนี้

% ────────────────────────────────────────────────────────────
\subsection{ผลการทดลองที่ 1: ความแม่นยำของการตรวจจับและนับยานพาหนะ}
\label{subsec:exp1}
% ────────────────────────────────────────────────────────────

ทำการเปรียบเทียบยอดนับจากระบบ AI Worker
กับค่าอ้างอิงจากการนับด้วยตาเปล่า (Ground Truth) ในชุดข้อมูลทดสอบ
โดยแยกประเมินตามประเภทยานพาหนะ
เพื่อทดสอบสมมติฐานด้านความยากในการตรวจจับรถจักรยานยนต์

\begin{table}[H]
  \centering
  \caption{ผลการทดสอบความแม่นยำของการตรวจจับและนับยานพาหนะแยกตามประเภท}
  \label{tab:exp1}
  \small
  \setlength{\tabcolsep}{5pt}
  \begin{tabular}{L{3.0cm} C{1.5cm} C{1.6cm} C{0.8cm} C{0.8cm} C{1.4cm} C{1.3cm} C{1.5cm} C{1.4cm}}
    \toprule
    \textbf{ประเภทยานพาหนะ} &
    \makecell{\textbf{Ground}\\\textbf{Truth}\\\textbf{(คัน)}} &
    \makecell{\textbf{System}\\\textbf{Count}\\\textbf{(TP)}} &
    \textbf{FP} &
    \textbf{FN} &
    \textbf{Precision} &
    \textbf{Recall} &
    \textbf{F1-Score} &
    \makecell{\textbf{MAE}\\\textbf{(คัน)}} \\
    \midrule
    รถยนต์ (Car)                       & 18          & 17          & 1 & 1 & 0.944 & 0.944 & 0.944 & 0.056 \\
    รถบรรทุก/รถบัส (Truck/Bus)         & 2           & 2           & 0 & 0 & 1.000 & 1.000 & 1.000 & 0.000 \\
    รถจักรยานยนต์ (Motorcycle)         & 8           & 7           & 1 & 1 & 0.875 & 0.875 & 0.875 & 0.125 \\
    \midrule
    \textbf{ภาพรวมทุกประเภท (Overall)} & \textbf{28} & \textbf{26} & \textbf{2} & \textbf{2}
        & \textbf{0.929} & \textbf{0.929} & \textbf{0.929} & \textbf{0.061} \\
    \bottomrule
  \end{tabular}
\end{table}

\noindent\textbf{วิเคราะห์ผล:}
ค่า Recall ของรถยนต์อยู่ที่ 0.944 (ผ่านเกณฑ์ $\geq 0.85$)
ค่า Recall ของรถจักรยานยนต์อยู่ที่ 0.875 (ผ่านเกณฑ์ $\geq 0.70$)
และค่า Precision ภาพรวมเท่ากับ 0.929 (ผ่านเกณฑ์ $\geq 0.85$)
เหตุผลที่รถจักรยานยนต์มี FN เกิดจากกรณีที่โมเดลจำแนกผิดเป็นชนิดจักรยาน (Bicycle)
ซึ่งถูกตั้งใจตัดออกตั้งแต่ชั้น AI Worker

% ────────────────────────────────────────────────────────────
\subsection{ผลการทดลองที่ 2: ความแม่นยำของการจำแนกสภาพการจราจร 4 สถานะ}
\label{subsec:exp2}
% ────────────────────────────────────────────────────────────

ประเมินการตัดสินสภาวะการจราจรในแต่ละหน้าต่างเวลา 10 วินาที
เทียบกับการกำกับสถานะจากวิดีโอจริง

\begin{table}[H]
  \centering
  \caption{Confusion Matrix การจำแนกสภาพการจราจร 4 สถานะ}
  \label{tab:confusion}
  \small
  \setlength{\tabcolsep}{4pt}
  \begin{tabular}{L{3.8cm} C{1.4cm} C{1.8cm} C{2.0cm} C{1.7cm} C{1.4cm} C{1.3cm}}
    \toprule
    \makecell[l]{\textbf{สภาพจริง} \\ $\backslash$\,\textbf{ผลประเมิน}} &
    \textbf{Normal} &
    \makecell{\textbf{High}\\\textbf{Density}} &
    \makecell{\textbf{Slow}\\\textbf{Moving}} &
    \textbf{Standstill} &
    \textbf{Precision} &
    \textbf{Recall} \\
    \midrule
    Normal ($O < 3.0$)
        & \textbf{12} & 0          & 0          & 0          & 1.000 & 1.000 \\
    High Density ($O \geq 3.0,\, V_c \geq 12.0$)
        & 1           & \textbf{4} & 0          & 0          & 1.000 & 0.800 \\
    Slow Moving ($O \geq 3.0,\, 2.0 \leq V_c < 12.0$)
        & 0           & 0          & \textbf{3} & 0          & 1.000 & 1.000 \\
    Standstill ($O \geq 3.0,\, V_c < 2.0$)
        & 0           & 0          & 0          & \textbf{2} & 1.000 & 1.000 \\
    \bottomrule
  \end{tabular}
\end{table}

\noindent\textbf{วิเคราะห์ผล:}
ความถูกต้องโดยรวม (Overall Accuracy) เท่ากับ \textbf{0.954 (21/22 windows)}
ซึ่งสูงกว่าเกณฑ์ที่กำหนดไว้ ($\geq 0.70$)
และที่สำคัญที่สุดคือ
\textbf{ไม่พบกรณีการตัดสินสภาวะหยุดนิ่ง (Standstill) เป็นการจราจรปกติ (Normal)}
ซึ่งสอดคล้องกับเกณฑ์การยอมรับขั้นเด็ดขาด (Must not occur)

% ────────────────────────────────────────────────────────────
\subsection{ผลการทดลองที่ 3: ผลของการแยกอัตราการไหลของรถจักรยานยนต์}
\label{subsec:exp3}
% ────────────────────────────────────────────────────────────

ทดสอบข้อเสนอหลักของโครงงาน
โดยเปรียบเทียบคำตัดสินสภาพการจราจรระหว่าง
\textbf{วิธีที่แยกอัตราการไหลของรถจักรยานยนต์ ($V_c$)}
กับ
\textbf{วิธีดั้งเดิมที่นับรวมยานพาหนะทุกประเภท ($V_{total}$)}

ข้อมูลทดสอบที่ใช้สาธิต (ช่วงวินาทีที่ 10--20):

\begin{lstlisting}[caption={ค่าพารามิเตอร์สภาพจราจรช่วงวินาทีที่ 10--20 ในการทดสอบ Experiment 3}]
ความหนาแน่นเฉลี่ย (O)          : 16.85 คัน/เฟรม  (รถยนต์จอดนิ่งเต็มโซน)
อัตราการไหลรถยนต์ (Vc)         : 0.0 คัน/นาที    (ไม่มีรถยนต์ผ่านเส้นนับ)
อัตราการไหลรถจักรยานยนต์ (Vmc) : 24.0 คัน/นาที  (รถจักรยานยนต์แทรกผ่านได้ 4 คัน)
\end{lstlisting}

\noindent\textbf{ผลการตัดสิน:}

\begin{itemize}
  \item \textbf{วิธีดั้งเดิม (นับรวม $V_{total} = 24.0$ คัน/นาที):}
        ระบบตัดสินเป็น \textbf{High Density}
        (จราจรหนาแน่นแต่ยังไหลได้ดี)
        ซึ่ง\textbf{ผิดจากความเป็นจริง} เพราะรถยนต์หยุดนิ่งสนิท

  \item \textbf{วิธีของโครงงาน (แยกเฉพาะ $V_c = 0.0$ คัน/นาที):}
        ระบบตัดสินเป็น \textbf{Standstill}
        (จราจรติดขัดหยุดนิ่ง)
        ซึ่ง\textbf{ถูกต้องตรงตามสภาพจริง}
\end{itemize}

\noindent\textbf{สรุป:}
มีหน้าต่างเวลาที่คำตัดสินเปลี่ยนไปอย่างมีนัยสำคัญตามเกณฑ์การยอมรับ
($\geq 1$ หน้าต่าง)
พิสูจน์ให้เห็นว่าการแยก $V_{mc}$ ออกจากเกณฑ์ตัดสิน
ช่วยป้องกันไม่ให้รถจักรยานยนต์กลบสัญญาณสภาวะรถติดรุนแรงได้จริง

% ────────────────────────────────────────────────────────────
\subsection{ผลการทดลองที่ 4: ความหน่วงของการเชื่อมโยงข้อมูล}
\label{subsec:exp4}
% ────────────────────────────────────────────────────────────

ทำการแทรกจุดประทับเวลา (Timestamping) 4 จุด
เพื่อวัดค่าความหน่วง End-to-End จาก AI Worker ถึงการแสดงผลบน Unity

\begin{table}[H]
  \centering
  \caption{ผลการวัดความหน่วงของระบบแบ่งตามช่วงการส่งผ่านข้อมูล}
  \label{tab:latency}
  \begin{tabular}{L{8.2cm} C{3.0cm} C{3.0cm}}
    \toprule
    \textbf{ช่วงการส่งผ่านข้อมูล (Processing Pipeline Segment)} &
    \textbf{Median Latency (ms)} &
    \textbf{P95 Latency (ms)} \\
    \midrule
    1.\ AI Worker Processing \& Event Generation              & 22.4          & 38.1 \\
    2.\ HTTP POST Ingest to Backend Validation                & 4.8           & 11.2 \\
    3.\ Backend WebSocket Broadcast (\texttt{spawn\_vehicle}) & 2.1           & 5.5  \\
    4.\ Unity WebSocket Receive to Object Pool Spawn          & 18.5          & 45.0 \\
    \midrule
    \textbf{รวมความหน่วงข้อมูลรายคัน (End-to-End Latency)}   & \textbf{47.8} & \textbf{99.8} \\
    \bottomrule
  \end{tabular}
\end{table}

\noindent\textbf{วิเคราะห์ผล:}
ความหน่วงรวมเฉลี่ย (Median Latency) ของเหตุการณ์รายคันเท่ากับ
\textbf{47.8 มิลลิวินาที (0.0478 วินาที)}
ซึ่งดีกว่าเกณฑ์ที่กำหนดไว้ไม่เกิน 1 วินาทีอย่างมาก
สำหรับคิวการสร้างใน Unity ไม่เกิดปัญหาสะสมคิวเนื่องจากอัตราการไหลของรถในวิดีโอ
ไม่เกินขีดจำกัดการดึงวัตถุจาก Object Pool (2 คัน/วินาที)

% ============================================================
\section{รายละเอียดและการทำงานของระบบต้นแบบ}
\label{sec:prototype}
% ============================================================

สถาปัตยกรรมระบบต้นแบบแฝดดิจิทัล (Digital Shadow level)
ประกอบด้วย 3 โมดูลหลักที่ทำงานร่วมกันอย่างเป็นระบบ
ดังแผนภาพในรูปที่~\ref{fig:pipeline}

\begin{figure}[H]
\begin{lstlisting}[language=,frame=single,basicstyle=\ttfamily\scriptsize,
                   caption={สถาปัตยกรรม End-to-End Pipeline ของระบบต้นแบบ Smart Flow},
                   label={fig:pipeline}]
+----------------------------------------------------------------------+
|                           PHYSICAL LAYER                             |
|           กล้อง CCTV สะพานลอย ถนนฉลองกรุง (3840x2160, 30fps)        |
+----------------------------------------------------------------------+
                                  | (Video Stream)
                                  v
+----------------------------------------------------------------------+
|               AI WORKER  (Python / YOLOv8s / ByteTrack)             |
|  - ตรวจจับ/ติดตามยานพาหนะ  (ImgSize=1280, Conf=0.35)                |
|  - คำนวณ Line Crossing & Density Polygon                             |
|  - สรุป Traffic State ทุก 10 วินาที (แยก Vc และ Vmc)                |
+----------------------------------------------------------------------+
         | HTTP POST  (/api/ingest  per vehicle, /api/traffic-state)
                                  v
+----------------------------------------------------------------------+
|            BACKEND SERVER  (Node.js / Express / Socket.IO)          |
|  - Data Contract Enforcement & HTTP 400 Validation                  |
|  - In-Memory Counter Accumulation                                   |
|  - WebSocket Broadcasting  ('spawn_vehicle', 'traffic_state')       |
+----------------------------------------------------------------------+
                       | WebSocket  (Socket.IO JSON)
                                  v
+----------------------------------------------------------------------+
|                  VIRTUAL LAYER  (Unity 3D / C# / WebGL)             |
|  - Event Listener & Queue Manager (FIFO)                            |
|  - Object Pooling Vehicle Spawning & Waypoint Navigation            |
|  - Real-time 3D Scene Environment Sync                              |
+----------------------------------------------------------------------+
\end{lstlisting}
\end{figure}

% ────────────────────────────────────────────────────────────
\subsection{ส่วนประสานผู้ใช้และการแสดงผลสามมิติ (Unity 3D Virtual Layer)}
\label{subsec:unity}
% ────────────────────────────────────────────────────────────

ส่วนแสดงผล 3D พัฒนาบน Unity เอนจินและปรับแต่งสำหรับ WebGL/PC
ประกอบด้วยองค์ประกอบหลัก 3 ส่วน

\begin{enumerate}
  \item \textbf{3D Virtual City View:}
        จำลองโครงสร้างทางกายภาพของถนนฉลองกรุงและสิ่งปลูกสร้างโดยรอบ
        สปอว์นโมเดล 3D ของรถยนต์ Bus Truck และ Motorcycle
        เคลื่อนที่ตามจุดนำทาง (Waypoints) บนเลนที่ถูกต้องแบบเรียลไทม์

  \item \textbf{Object Pooling System:}
        สร้าง Prefab ยานพาหนะสำรองไว้ในหน่วยความจำล่วงหน้า
        เมื่อได้รับอีเวนต์ \texttt{spawn\_vehicle}
        จะดึงวัตถุมาเปิดใช้งาน (Enable)
        และคืนค่าเข้าคิว (Disable) เมื่อวิ่งสุดเส้นทาง
        ช่วยลดอาการกระตุกของ Frame Rate บน Unity

  \item \textbf{Developer Monitoring Dashboard:}
        หน้าเว็บสำหรับผู้พัฒนาที่รับข้อมูลผ่าน WebSocket ช่องทาง
        \texttt{spawn} และ \texttt{traffic\_state}
        แสดงสถานะตัวนับสะสม (Counter Statistics)
        และผลตัดสินสภาพการจราจรในรูปแบบตัวเลขและกราฟแบบ Real-time
\end{enumerate}

% ────────────────────────────────────────────────────────────
\subsection{สภาพการณ์การใช้งานระบบ (Use Case Scenarios)}
\label{subsec:scenarios}
% ────────────────────────────────────────────────────────────

ระบบรองรับการทำงานในสถานการณ์จริง 4 รูปแบบหลัก

\begin{itemize}
  \item \textbf{Scenario 1 --- Normal Traffic Flow Sync:}
        ยานพาหนะข้ามเส้นนับ
        AI Worker ส่งเหตุการณ์รายคันผ่าน HTTP POST
        Backend ตรวจสอบผ่าน Data Contract แล้วส่งต่อผ่าน WebSocket
        ให้ Unity สปอว์นโมเดลรถวิ่งด้วยความเร็วปกติทันที

  \item \textbf{Scenario 2 --- Standstill Detection during Mixed Traffic:}
        เกิดสภาวะรถยนต์ติดสะสมหยุดนิ่ง แต่รถจักรยานยนต์ยังแทรกตัวผ่านได้
        AI Worker สรุปค่า $O = 16.85$ และ $V_c = 0.0$
        ส่งคำตัดสิน \texttt{standstill} ไปยัง Backend
        และกระจายไปปรับสภาวะฉาก 3D

  \item \textbf{Scenario 3 --- Invalid Data Handling (Data Contract Gate):}
        หากโมดูลต้นทางส่งข้อมูลที่มีโครงสร้างผิดพลาด
        เช่น \texttt{confidence} ไม่อยู่ในช่วง 0.0--1.0 หรือเวลาติดลบ
        Backend จะปฏิเสธคำขอด้วยรหัส \texttt{HTTP 400 Bad Request} ทันที
        ป้องกันไม่ให้ข้อมูลขยะหลุดไปสร้างความผิดพลาดในส่วนแสดงผล 3D

  \item \textbf{Scenario 4 --- Track Loss \& Auto-Recovery:}
        หากยานพาหนะถูกบดบังชั่วขณะ
        อัลกอริทึม ByteTrack จะรักษาสถานะ Tracklet ไว้ในคิว Lost Track
        เป็นเวลา 30 เฟรม ก่อนจะจับคู่คืนรหัส ID เดิมเมื่อวัตถุปรากฏขึ้นอีกครั้ง
        ป้องกันการสร้างยานพาหนะซ้ำซ้อนในโลกเสมือน
\end{itemize}

% ============================================================
\section{สรุปการประเมินผลเทียบกับเกณฑ์การยอมรับ}
\label{sec:summary}
% ============================================================

จากการทดสอบระบบตามตัวชี้วัดทั้งหมดที่กำหนดไว้ในบทที่ 3
สามารถสรุปผลการประเมินเปรียบเทียบได้ดังตารางที่~\ref{tab:acceptance}

\begin{table}[H]
  \centering
  \caption{สรุปผลการประเมินเทียบกับเกณฑ์การยอมรับทั้งหมด}
  \label{tab:acceptance}
  \small
  \setlength{\tabcolsep}{5pt}
  \begin{tabular}{L{4.5cm} C{3.2cm} C{3.0cm} C{2.5cm}}
    \toprule
    \textbf{การทดลอง / ตัวชี้วัด} &
    \makecell{\textbf{เกณฑ์การยอมรับ}\\\textbf{(Target)}} &
    \makecell{\textbf{ผลที่ได้}\\\textbf{(Actual Result)}} &
    \textbf{การประเมิน} \\
    \midrule
    Exp 1: Recall ของรถยนต์ (Car)
        & $\geq 0.85$
        & \textbf{0.944}
        & \textbf{ผ่านเกณฑ์} \\
    Exp 1: Recall ของรถจักรยานยนต์ (Motorcycle)
        & $\geq 0.70$
        & \textbf{0.875}
        & \textbf{ผ่านเกณฑ์} \\
    Exp 1: Precision รวมทุกประเภท
        & $\geq 0.85$
        & \textbf{0.929}
        & \textbf{ผ่านเกณฑ์} \\
    Exp 2: ความถูกต้องรวมการจำแนกสถานะ
        & $\geq 0.70$
        & \textbf{0.954}
        & \textbf{ผ่านเกณฑ์} \\
    Exp 2: ตัดสิน Standstill เป็น Normal
        & ต้องไม่เกิดขึ้นเลย
        & \textbf{0 กรณี}
        & \textbf{ผ่านเกณฑ์} \\
    Exp 3: หน้าต่างที่คำตัดสินเปลี่ยน
        & $\geq 1$ หน้าต่าง
        & \textbf{1 หน้าต่าง}
        & \textbf{ผ่านเกณฑ์} \\
    Exp 4: ความหน่วงข้อมูลรายคัน (Median)
        & $\leq 1.0$ วินาที
        & \textbf{0.0478 วินาที}
        & \textbf{ผ่านเกณฑ์} \\
    \bottomrule
  \end{tabular}
\end{table}

ผลการทดลองในบทนี้ยืนยันว่าระบบต้นแบบแฝดดิจิทัลที่พัฒนาขึ้น
สามารถประมวลผลวิดีโอ แยกประเภทยานพาหนะ
ตัดสินสภาพการจราจรแบบผสมในประเทศไทยได้อย่างแม่นยำ
และส่งต่อข้อมูลไปจำลองภาพ 3D แบบเรียลไทม์ได้ตรงตามเป้าหมายของโครงงานทุกประการ

% ============================================================
\end{document}
% ============================================================

% ============================================================

\documentclass[12pt,a4paper]{report}

% ── Thai / Font ─────────────────────────────────────────────
\usepackage{fontspec}
\setmainfont{TH Sarabun New}
\setsansfont{TH Sarabun New}
\setmonofont[Scale=0.82]{Courier New}

% ── Layout ──────────────────────────────────────────────────
\usepackage[
  a4paper,
  left=3.5cm, right=2.5cm,
  top=3.0cm,  bottom=3.0cm
]{geometry}
\usepackage{setspace}
\onehalfspacing
\setlength{\parindent}{1.5em}
\setlength{\parskip}{0.3em}

% ── Math ────────────────────────────────────────────────────
\usepackage{amsmath}
\usepackage{amssymb}

% ── Tables ──────────────────────────────────────────────────
\usepackage{booktabs}
\usepackage{array}
\usepackage{multirow}
\usepackage{makecell}
\usepackage{longtable}
\usepackage{tabularx}
\renewcommand{\arraystretch}{1.5}

% ── Figures ─────────────────────────────────────────────────
\usepackage{graphicx}
\usepackage{float}
\usepackage{caption}
\captionsetup{font=small, labelfont=bf}

% ── Code blocks ─────────────────────────────────────────────
\usepackage{listings}
\usepackage{xcolor}

\definecolor{lstbg}{RGB}{248,248,248}
\definecolor{lstframe}{RGB}{200,200,200}
\definecolor{lstcmt}{RGB}{60,140,60}
\definecolor{lstkw}{RGB}{0,100,160}

\lstset{
  backgroundcolor = \color{lstbg},
  frame           = single,
  rulecolor       = \color{lstframe},
  basicstyle      = \ttfamily\small,
  commentstyle    = \color{lstcmt},
  keywordstyle    = \color{lstkw}\bfseries,
  breaklines      = true,
  numbers         = none,
  captionpos      = b,
  xleftmargin     = 0.8em,
  xrightmargin    = 0.8em,
  aboveskip       = 0.8em,
  belowskip       = 0.6em,
  keepspaces      = true,
  showstringspaces= false,
}

% ── Hyperlinks ───────────────────────────────────────────────
\usepackage[hidelinks]{hyperref}

% ── Custom column types ──────────────────────────────────────
\newcolumntype{C}[1]{>{\centering\arraybackslash}p{#1}}
\newcolumntype{L}[1]{>{\raggedright\arraybackslash}p{#1}}
\newcolumntype{R}[1]{>{\raggedleft\arraybackslash}p{#1}}

% ============================================================
\begin{document}
% ============================================================

% ตั้ง counter ให้บทนี้เป็นบทที่ 4 (หากใช้ \input ในเอกสารหลักให้ลบบรรทัดนี้ออก)
\setcounter{chapter}{3}

\chapter{ผลการทดลองเบื้องต้นและระบบต้นแบบ}
\label{chap:results}

% ============================================================
\section{ผลการทดลองเบื้องต้นและการเลือกองค์ประกอบระบบ}
\label{sec:phase1}
% ============================================================

ในส่วนนี้เป็นการทดลองเบื้องต้นเพื่อคัดเลือกและปรับแต่งค่าพารามิเตอร์ขององค์ประกอบหลัก
ใน Pipeline การประมวลผลวิดีโอ (Video Processing \& AI Pipeline)
เพื่อให้ระบบมีความเหมาะสมและสมดุลที่สุดระหว่างความแม่นยำในการตรวจจับยานพาหนะขนาดเล็ก
กับการทำงานแบบเรียลไทม์

% ────────────────────────────────────────────────────────────
\subsection{การทดสอบขนาดภาพนำเข้าและขีดแบ่งความเชื่อมั่น}
\label{subsec:input-size}
% ────────────────────────────────────────────────────────────

ในการประมวลผลวิดีโอความละเอียด $3840 \times 2160$ พิกเซล จากกล้องมุมสูงบนสะพานลอยถนนฉลองกรุง
พบว่ารถจักรยานยนต์ซึ่งเป็นองค์ประกอบหลักของกระแสจราจรแบบผสม (Mixed Traffic)
ปรากฏเป็นวัตถุขนาดเล็กเมื่ออยู่ไกลจากกล้อง
จึงได้ทำการทดสอบ 2 พารามิเตอร์ที่มีผลต่อการตรวจจับยานพาหนะขนาดเล็กโดยตรง ดังนี้

\paragraph{การทดสอบ Input Image Size}
เปรียบเทียบ imgsz = 640 (ค่ามาตรฐาน YOLOv8) กับ imgsz = 1{,}280
บนวิดีโอตัวอย่าง 525 เฟรม (1920$\times$1080, 25 fps, 21 วินาที)
โดยคงพารามิเตอร์อื่นไว้คงที่ (conf = 0.35, ByteTrack)

\begin{table}[H]
  \centering
  \caption{ผลเปรียบเทียบ Input Image Size ต่อจำนวนยานพาหนะที่นับได้}
  \label{tab:imgsz}
  \small
  \begin{tabular}{L{3.5cm} C{1.8cm} C{2.2cm} C{2.0cm} C{2.0cm} C{2.0cm}}
    \toprule
    \textbf{imgsz} & \textbf{Car} & \textbf{Motorcycle} & \textbf{Truck/Bus} & \textbf{Total} & \textbf{FPS} \\
    \midrule
    640 (ค่ามาตรฐาน) & 5 & \textbf{0} & 0 &  5 & 22.9 \\
    \textbf{1280 (เลือก)} & 3 & \textbf{8} & 0 & 11 & \textbf{18.7} \\
    \bottomrule
  \end{tabular}
\end{table}

\noindent\textbf{วิเคราะห์ผล:}
imgsz = 640 ตรวจจับรถจักรยานยนต์ได้ \textbf{0 คัน} (Recall = 0.000)
เนื่องจากภาพถูกย่อลง 3 เท่าทำให้พิกเซลของรถจักรยานยนต์เล็กเกินขีดจำกัดของโมเดล
imgsz = 1{,}280 ตรวจจับได้ครบ \textbf{8 คัน}
โดยใช้เวลาเพิ่มขึ้นเพียง 5.1 วินาที (จาก 22.9 เป็น 18.7 FPS)
ซึ่งยังอยู่ในระดับที่เหมาะสมสำหรับการทำงานแบบ Real-time

\paragraph{การทดสอบ Confidence Threshold}
ทำการทดสอบค่าขีดแบ่ง 4 ระดับ โดยตรึง imgsz = 1{,}280 ไว้

\begin{table}[H]
  \centering
  \caption{ผลเปรียบเทียบ Confidence Threshold ต่อจำนวน Motorcycle ที่นับได้ (imgsz = 1{,}280)}
  \label{tab:conf}
  \small
  \begin{tabular}{C{3.2cm} C{1.8cm} C{2.5cm} C{2.0cm} C{3.5cm}}
    \toprule
    \textbf{Threshold} & \textbf{Car} & \textbf{Motorcycle} & \textbf{Total} & \textbf{FP (conf $<$ 0.40)} \\
    \midrule
    0.60 (ค่ามาตรฐาน) & 3 & 0            &  3 & 0 \\
    0.50               & 3 & 2            &  5 & 0 \\
    0.40               & 3 & 6            &  9 & 0 \\
    \textbf{0.35 (เลือก)} & 3 & \textbf{8} & \textbf{11} & \textbf{0} \\
    \bottomrule
  \end{tabular}
\end{table}

\noindent\textbf{วิเคราะห์ผล:}
จากการวัดพบว่า Confidence ของรถจักรยานยนต์ ณ เวลาที่ข้ามเส้นนับอยู่ในช่วง
$0.400 \leq \text{conf} \leq 0.549$ (ค่าเฉลี่ย 0.466)
ต่ำกว่ารถยนต์ที่อยู่ในช่วง $0.678 \leq \text{conf} \leq 0.834$ อย่างเป็นระบบ
การใช้ Threshold = 0.60 จึงตัดรถจักรยานยนต์ออกทั้งหมด
ระบบกำหนดค่าขีดแบ่งไว้ที่ \textbf{0.35} ซึ่งคงข้อมูลรถจักรยานยนต์ครบ 8 คัน
โดยไม่มี False Positive เพิ่มเลย เนื่องจากกระบวนการ Vote ของ ByteTrack
กรอง Track ที่ไม่น่าเชื่อถือออกก่อนการนับ

% ────────────────────────────────────────────────────────────
\subsection{การเปรียบเทียบและคัดเลือกอัลกอริทึมติดตามวัตถุ}
\label{subsec:tracking}
% ────────────────────────────────────────────────────────────

ในการทดสอบอัลกอริทึมติดตามวัตถุหลายเป้าหมาย (Multi-Object Tracking)
เพื่อป้องกันการนับยานพาหนะซ้ำ ได้ทำการเปรียบเทียบ 2 อัลกอริทึมหลัก
บนเงื่อนไขเดียวกัน (imgsz = 1{,}280, conf = 0.35, วิดีโอตัวอย่าง 525 เฟรม)

\begin{table}[H]
  \centering
  \caption{ผลเปรียบเทียบ Tracking Algorithm (ByteTrack vs.\ BoT-SORT)}
  \label{tab:tracker}
  \small
  \begin{tabular}{L{6.5cm} C{3.0cm} C{3.0cm}}
    \toprule
    \textbf{ตัวชี้วัด} & \textbf{ByteTrack} & \textbf{BoT-SORT} \\
    \midrule
    Crossing events รวม (คัน)              & 11            & 11       \\
    \quad --- Car                           & 3             & 3        \\
    \quad --- Motorcycle                    & 8             & 8        \\
    ความเร็วประมวลผล (FPS)                  & \textbf{16.4} & 11.5     \\
    Max Motorcycle Track ID (ID Churn)      & 157           & 150      \\
    Anomalies (การนับที่ถูกปฏิเสธ)          & 0             & 0        \\
    Crossing timestamps ตรงกัน             & \checkmark    & \checkmark \\
    \bottomrule
  \end{tabular}
\end{table}

\noindent\textbf{วิเคราะห์ผล:}
ทั้งสองอัลกอริทึมนับยานพาหนะได้จำนวนเท่ากัน (11 คัน)
และตรวจพบยานพาหนะในเวลาเดียวกันทุกรายการ
แสดงให้เห็นว่า Counting Quality ไม่แตกต่างกัน
ByteTrack มีความเร็วสูงกว่า BoT-SORT ถึง \textbf{30\%} (16.4 vs.\ 11.5 FPS)
ในขณะที่ ID Churn ของ Motorcycle ที่ BoT-SORT ต่ำกว่าเพียง 4\%
(Max Track ID 150 vs.\ 157) ซึ่งไม่คุ้มค่ากับ FPS ที่สูญเสียไป
ระบบจึงคัดเลือก \textbf{ByteTrack} เป็นค่าตั้งต้นสำหรับ AI Worker

\medskip
\noindent\textbf{หมายเหตุเกี่ยวกับ ID Churn:}
ค่า Max Motorcycle Track ID สะท้อนจำนวน Track ทั้งหมดที่ถูกสร้างขึ้นตลอดคลิป
(ByteTrack สร้าง 157 Track สำหรับรถจักรยานยนต์ที่ข้ามเส้น 8 คัน = อัตราส่วน 19.6:1)
ค่าที่สูงนี้เป็นปกติสำหรับจราจรแบบผสม เนื่องจาก Tracker สร้าง Track ชั่วคราว
สำหรับยานพาหนะที่เข้ามาในโซนแต่ยังไม่ถึงเส้นนับ
ByteTrack จัดการ Track เหล่านี้ด้วย Two-stage Data Association
ซึ่งฟื้นฟู Track ที่หายไปชั่วคราวได้ โดยไม่เพิ่มจำนวน Crossing Event ที่ผิดพลาด

% ────────────────────────────────────────────────────────────
\subsection{การกำหนดขนาดหน้าต่างเวลาประเมินสภาพจราจร}
\label{subsec:time-window}
% ────────────────────────────────────────────────────────────

ระบบทำการทดสอบความยาวหน้าต่างเวลา ($T$) ในการคำนวณอัตราการไหล ($V$)
และค่าความหนาแน่น ($O$)
พบว่าหน้าต่างเวลาที่ \textbf{10 วินาที} ให้ค่าเฉลี่ยสถิติการข้ามเส้น
ในระดับหลักหน่วยถึงหลักสิบ
ซึ่งไม่ผันผวนรวดเร็วเกินไปจนระบบแกว่ง (เมื่อเทียบกับ 5 วินาที)
และไม่ล่าช้าเกินไปในการตอบสนองต่อสภาวะรถติด (เมื่อเทียบกับ 30 วินาที)

% ============================================================
\section{ผลการทดสอบประสิทธิภาพของระบบและการจำแนกสถานะการจราจร}
\label{sec:phase2}
% ============================================================

ผลการทดลองทั้ง 4 ชุดตามที่ออกแบบไว้ในแผนงาน
เพื่อประเมินความถูกต้องของการนับ การจำแนกสภาวะ
ผลของการแยกอัตราการไหล และความหน่วงของข้อมูล สรุปได้ดังนี้

% ────────────────────────────────────────────────────────────
\subsection{ผลการทดลองที่ 1: ความแม่นยำของการตรวจจับและนับยานพาหนะ}
\label{subsec:exp1}
% ────────────────────────────────────────────────────────────

ทำการเปรียบเทียบยอดนับจากระบบ AI Worker
กับค่าอ้างอิงจากการนับด้วยตาเปล่า (Ground Truth) ในชุดข้อมูลทดสอบ
โดยแยกประเมินตามประเภทยานพาหนะ
เพื่อทดสอบสมมติฐานด้านความยากในการตรวจจับรถจักรยานยนต์

\begin{table}[H]
  \centering
  \caption{ผลการทดสอบความแม่นยำของการตรวจจับและนับยานพาหนะแยกตามประเภท}
  \label{tab:exp1}
  \small
  \setlength{\tabcolsep}{5pt}
  \begin{tabular}{L{3.0cm} C{1.5cm} C{1.6cm} C{0.8cm} C{0.8cm} C{1.4cm} C{1.3cm} C{1.5cm} C{1.4cm}}
    \toprule
    \textbf{ประเภทยานพาหนะ} &
    \makecell{\textbf{Ground}\\\textbf{Truth}\\\textbf{(คัน)}} &
    \makecell{\textbf{System}\\\textbf{Count}\\\textbf{(TP)}} &
    \textbf{FP} &
    \textbf{FN} &
    \textbf{Precision} &
    \textbf{Recall} &
    \textbf{F1-Score} &
    \makecell{\textbf{MAE}\\\textbf{(คัน)}} \\
    \midrule
    รถยนต์ (Car)                       & 18          & 17          & 1 & 1 & 0.944 & 0.944 & 0.944 & 0.056 \\
    รถบรรทุก/รถบัส (Truck/Bus)         & 2           & 2           & 0 & 0 & 1.000 & 1.000 & 1.000 & 0.000 \\
    รถจักรยานยนต์ (Motorcycle)         & 8           & 7           & 1 & 1 & 0.875 & 0.875 & 0.875 & 0.125 \\
    \midrule
    \textbf{ภาพรวมทุกประเภท (Overall)} & \textbf{28} & \textbf{26} & \textbf{2} & \textbf{2}
        & \textbf{0.929} & \textbf{0.929} & \textbf{0.929} & \textbf{0.061} \\
    \bottomrule
  \end{tabular}
\end{table}

\noindent\textbf{วิเคราะห์ผล:}
ค่า Recall ของรถยนต์อยู่ที่ 0.944 (ผ่านเกณฑ์ $\geq 0.85$)
ค่า Recall ของรถจักรยานยนต์อยู่ที่ 0.875 (ผ่านเกณฑ์ $\geq 0.70$)
และค่า Precision ภาพรวมเท่ากับ 0.929 (ผ่านเกณฑ์ $\geq 0.85$)
เหตุผลที่รถจักรยานยนต์มี FN เกิดจากกรณีที่โมเดลจำแนกผิดเป็นชนิดจักรยาน (Bicycle)
ซึ่งถูกตั้งใจตัดออกตั้งแต่ชั้น AI Worker

% ────────────────────────────────────────────────────────────
\subsection{ผลการทดลองที่ 2: ความแม่นยำของการจำแนกสภาพการจราจร 4 สถานะ}
\label{subsec:exp2}
% ────────────────────────────────────────────────────────────

ประเมินการตัดสินสภาวะการจราจรในแต่ละหน้าต่างเวลา 10 วินาที
เทียบกับการกำกับสถานะจากวิดีโอจริง

\begin{table}[H]
  \centering
  \caption{Confusion Matrix การจำแนกสภาพการจราจร 4 สถานะ}
  \label{tab:confusion}
  \small
  \setlength{\tabcolsep}{4pt}
  \begin{tabular}{L{3.8cm} C{1.4cm} C{1.8cm} C{2.0cm} C{1.7cm} C{1.4cm} C{1.3cm}}
    \toprule
    \makecell[l]{\textbf{สภาพจริง} \\ $\backslash$\,\textbf{ผลประเมิน}} &
    \textbf{Normal} &
    \makecell{\textbf{High}\\\textbf{Density}} &
    \makecell{\textbf{Slow}\\\textbf{Moving}} &
    \textbf{Standstill} &
    \textbf{Precision} &
    \textbf{Recall} \\
    \midrule
    Normal ($O < 3.0$)
        & \textbf{12} & 0          & 0          & 0          & 1.000 & 1.000 \\
    High Density ($O \geq 3.0,\, V_c \geq 12.0$)
        & 1           & \textbf{4} & 0          & 0          & 1.000 & 0.800 \\
    Slow Moving ($O \geq 3.0,\, 2.0 \leq V_c < 12.0$)
        & 0           & 0          & \textbf{3} & 0          & 1.000 & 1.000 \\
    Standstill ($O \geq 3.0,\, V_c < 2.0$)
        & 0           & 0          & 0          & \textbf{2} & 1.000 & 1.000 \\
    \bottomrule
  \end{tabular}
\end{table}

\noindent\textbf{วิเคราะห์ผล:}
ความถูกต้องโดยรวม (Overall Accuracy) เท่ากับ \textbf{0.954 (21/22 windows)}
ซึ่งสูงกว่าเกณฑ์ที่กำหนดไว้ ($\geq 0.70$)
และที่สำคัญที่สุดคือ
\textbf{ไม่พบกรณีการตัดสินสภาวะหยุดนิ่ง (Standstill) เป็นการจราจรปกติ (Normal)}
ซึ่งสอดคล้องกับเกณฑ์การยอมรับขั้นเด็ดขาด (Must not occur)

% ────────────────────────────────────────────────────────────
\subsection{ผลการทดลองที่ 3: ผลของการแยกอัตราการไหลของรถจักรยานยนต์}
\label{subsec:exp3}
% ────────────────────────────────────────────────────────────

ทดสอบข้อเสนอหลักของโครงงาน
โดยเปรียบเทียบคำตัดสินสภาพการจราจรระหว่าง
\textbf{วิธีที่แยกอัตราการไหลของรถจักรยานยนต์ ($V_c$)}
กับ
\textbf{วิธีดั้งเดิมที่นับรวมยานพาหนะทุกประเภท ($V_{total}$)}

ข้อมูลทดสอบที่ใช้สาธิต (ช่วงวินาทีที่ 10--20):

\begin{lstlisting}[caption={ค่าพารามิเตอร์สภาพจราจรช่วงวินาทีที่ 10--20 ในการทดสอบ Experiment 3}]
ความหนาแน่นเฉลี่ย (O)          : 16.85 คัน/เฟรม  (รถยนต์จอดนิ่งเต็มโซน)
อัตราการไหลรถยนต์ (Vc)         : 0.0 คัน/นาที    (ไม่มีรถยนต์ผ่านเส้นนับ)
อัตราการไหลรถจักรยานยนต์ (Vmc) : 24.0 คัน/นาที  (รถจักรยานยนต์แทรกผ่านได้ 4 คัน)
\end{lstlisting}

\noindent\textbf{ผลการตัดสิน:}

\begin{itemize}
  \item \textbf{วิธีดั้งเดิม (นับรวม $V_{total} = 24.0$ คัน/นาที):}
        ระบบตัดสินเป็น \textbf{High Density}
        (จราจรหนาแน่นแต่ยังไหลได้ดี)
        ซึ่ง\textbf{ผิดจากความเป็นจริง} เพราะรถยนต์หยุดนิ่งสนิท

  \item \textbf{วิธีของโครงงาน (แยกเฉพาะ $V_c = 0.0$ คัน/นาที):}
        ระบบตัดสินเป็น \textbf{Standstill}
        (จราจรติดขัดหยุดนิ่ง)
        ซึ่ง\textbf{ถูกต้องตรงตามสภาพจริง}
\end{itemize}

\noindent\textbf{สรุป:}
มีหน้าต่างเวลาที่คำตัดสินเปลี่ยนไปอย่างมีนัยสำคัญตามเกณฑ์การยอมรับ
($\geq 1$ หน้าต่าง)
พิสูจน์ให้เห็นว่าการแยก $V_{mc}$ ออกจากเกณฑ์ตัดสิน
ช่วยป้องกันไม่ให้รถจักรยานยนต์กลบสัญญาณสภาวะรถติดรุนแรงได้จริง

% ────────────────────────────────────────────────────────────
\subsection{ผลการทดลองที่ 4: ความหน่วงของการเชื่อมโยงข้อมูล}
\label{subsec:exp4}
% ────────────────────────────────────────────────────────────

ทำการแทรกจุดประทับเวลา (Timestamping) 4 จุด
เพื่อวัดค่าความหน่วง End-to-End จาก AI Worker ถึงการแสดงผลบน Unity

\begin{table}[H]
  \centering
  \caption{ผลการวัดความหน่วงของระบบแบ่งตามช่วงการส่งผ่านข้อมูล}
  \label{tab:latency}
  \begin{tabular}{L{8.2cm} C{3.0cm} C{3.0cm}}
    \toprule
    \textbf{ช่วงการส่งผ่านข้อมูล (Processing Pipeline Segment)} &
    \textbf{Median Latency (ms)} &
    \textbf{P95 Latency (ms)} \\
    \midrule
    1.\ AI Worker Processing \& Event Generation              & 22.4          & 38.1 \\
    2.\ HTTP POST Ingest to Backend Validation                & 4.8           & 11.2 \\
    3.\ Backend WebSocket Broadcast (\texttt{spawn\_vehicle}) & 2.1           & 5.5  \\
    4.\ Unity WebSocket Receive to Object Pool Spawn          & 18.5          & 45.0 \\
    \midrule
    \textbf{รวมความหน่วงข้อมูลรายคัน (End-to-End Latency)}   & \textbf{47.8} & \textbf{99.8} \\
    \bottomrule
  \end{tabular}
\end{table}

\noindent\textbf{วิเคราะห์ผล:}
ความหน่วงรวมเฉลี่ย (Median Latency) ของเหตุการณ์รายคันเท่ากับ
\textbf{47.8 มิลลิวินาที (0.0478 วินาที)}
ซึ่งดีกว่าเกณฑ์ที่กำหนดไว้ไม่เกิน 1 วินาทีอย่างมาก
สำหรับคิวการสร้างใน Unity ไม่เกิดปัญหาสะสมคิวเนื่องจากอัตราการไหลของรถในวิดีโอ
ไม่เกินขีดจำกัดการดึงวัตถุจาก Object Pool (2 คัน/วินาที)

% ============================================================
\section{รายละเอียดและการทำงานของระบบต้นแบบ}
\label{sec:prototype}
% ============================================================

สถาปัตยกรรมระบบต้นแบบแฝดดิจิทัล (Digital Shadow level)
ประกอบด้วย 3 โมดูลหลักที่ทำงานร่วมกันอย่างเป็นระบบ
ดังแผนภาพในรูปที่~\ref{fig:pipeline}

\begin{figure}[H]
\begin{lstlisting}[language=,frame=single,basicstyle=\ttfamily\scriptsize,
                   caption={สถาปัตยกรรม End-to-End Pipeline ของระบบต้นแบบ Smart Flow},
                   label={fig:pipeline}]
+----------------------------------------------------------------------+
|                           PHYSICAL LAYER                             |
|           กล้อง CCTV สะพานลอย ถนนฉลองกรุง (3840x2160, 30fps)        |
+----------------------------------------------------------------------+
                                  | (Video Stream)
                                  v
+----------------------------------------------------------------------+
|               AI WORKER  (Python / YOLOv8s / ByteTrack)             |
|  - ตรวจจับ/ติดตามยานพาหนะ  (ImgSize=1280, Conf=0.35)                |
|  - คำนวณ Line Crossing & Density Polygon                             |
|  - สรุป Traffic State ทุก 10 วินาที (แยก Vc และ Vmc)                |
+----------------------------------------------------------------------+
         | HTTP POST  (/api/ingest  per vehicle, /api/traffic-state)
                                  v
+----------------------------------------------------------------------+
|            BACKEND SERVER  (Node.js / Express / Socket.IO)          |
|  - Data Contract Enforcement & HTTP 400 Validation                  |
|  - In-Memory Counter Accumulation                                   |
|  - WebSocket Broadcasting  ('spawn_vehicle', 'traffic_state')       |
+----------------------------------------------------------------------+
                       | WebSocket  (Socket.IO JSON)
                                  v
+----------------------------------------------------------------------+
|                  VIRTUAL LAYER  (Unity 3D / C# / WebGL)             |
|  - Event Listener & Queue Manager (FIFO)                            |
|  - Object Pooling Vehicle Spawning & Waypoint Navigation            |
|  - Real-time 3D Scene Environment Sync                              |
+----------------------------------------------------------------------+
\end{lstlisting}
\end{figure}

% ────────────────────────────────────────────────────────────
\subsection{ส่วนประสานผู้ใช้และการแสดงผลสามมิติ (Unity 3D Virtual Layer)}
\label{subsec:unity}
% ────────────────────────────────────────────────────────────

ส่วนแสดงผล 3D พัฒนาบน Unity เอนจินและปรับแต่งสำหรับ WebGL/PC
ประกอบด้วยองค์ประกอบหลัก 3 ส่วน

\begin{enumerate}
  \item \textbf{3D Virtual City View:}
        จำลองโครงสร้างทางกายภาพของถนนฉลองกรุงและสิ่งปลูกสร้างโดยรอบ
        สปอว์นโมเดล 3D ของรถยนต์ Bus Truck และ Motorcycle
        เคลื่อนที่ตามจุดนำทาง (Waypoints) บนเลนที่ถูกต้องแบบเรียลไทม์

  \item \textbf{Object Pooling System:}
        สร้าง Prefab ยานพาหนะสำรองไว้ในหน่วยความจำล่วงหน้า
        เมื่อได้รับอีเวนต์ \texttt{spawn\_vehicle}
        จะดึงวัตถุมาเปิดใช้งาน (Enable)
        และคืนค่าเข้าคิว (Disable) เมื่อวิ่งสุดเส้นทาง
        ช่วยลดอาการกระตุกของ Frame Rate บน Unity

  \item \textbf{Developer Monitoring Dashboard:}
        หน้าเว็บสำหรับผู้พัฒนาที่รับข้อมูลผ่าน WebSocket ช่องทาง
        \texttt{spawn} และ \texttt{traffic\_state}
        แสดงสถานะตัวนับสะสม (Counter Statistics)
        และผลตัดสินสภาพการจราจรในรูปแบบตัวเลขและกราฟแบบ Real-time
\end{enumerate}

% ────────────────────────────────────────────────────────────
\subsection{สภาพการณ์การใช้งานระบบ (Use Case Scenarios)}
\label{subsec:scenarios}
% ────────────────────────────────────────────────────────────

ระบบรองรับการทำงานในสถานการณ์จริง 4 รูปแบบหลัก

\begin{itemize}
  \item \textbf{Scenario 1 --- Normal Traffic Flow Sync:}
        ยานพาหนะข้ามเส้นนับ
        AI Worker ส่งเหตุการณ์รายคันผ่าน HTTP POST
        Backend ตรวจสอบผ่าน Data Contract แล้วส่งต่อผ่าน WebSocket
        ให้ Unity สปอว์นโมเดลรถวิ่งด้วยความเร็วปกติทันที

  \item \textbf{Scenario 2 --- Standstill Detection during Mixed Traffic:}
        เกิดสภาวะรถยนต์ติดสะสมหยุดนิ่ง แต่รถจักรยานยนต์ยังแทรกตัวผ่านได้
        AI Worker สรุปค่า $O = 16.85$ และ $V_c = 0.0$
        ส่งคำตัดสิน \texttt{standstill} ไปยัง Backend
        และกระจายไปปรับสภาวะฉาก 3D

  \item \textbf{Scenario 3 --- Invalid Data Handling (Data Contract Gate):}
        หากโมดูลต้นทางส่งข้อมูลที่มีโครงสร้างผิดพลาด
        เช่น \texttt{confidence} ไม่อยู่ในช่วง 0.0--1.0 หรือเวลาติดลบ
        Backend จะปฏิเสธคำขอด้วยรหัส \texttt{HTTP 400 Bad Request} ทันที
        ป้องกันไม่ให้ข้อมูลขยะหลุดไปสร้างความผิดพลาดในส่วนแสดงผล 3D

  \item \textbf{Scenario 4 --- Track Loss \& Auto-Recovery:}
        หากยานพาหนะถูกบดบังชั่วขณะ
        อัลกอริทึม ByteTrack จะรักษาสถานะ Tracklet ไว้ในคิว Lost Track
        เป็นเวลา 30 เฟรม ก่อนจะจับคู่คืนรหัส ID เดิมเมื่อวัตถุปรากฏขึ้นอีกครั้ง
        ป้องกันการสร้างยานพาหนะซ้ำซ้อนในโลกเสมือน
\end{itemize}

% ============================================================
\section{สรุปการประเมินผลเทียบกับเกณฑ์การยอมรับ}
\label{sec:summary}
% ============================================================

จากการทดสอบระบบตามตัวชี้วัดทั้งหมดที่กำหนดไว้ในบทที่ 3
สามารถสรุปผลการประเมินเปรียบเทียบได้ดังตารางที่~\ref{tab:acceptance}

\begin{table}[H]
  \centering
  \caption{สรุปผลการประเมินเทียบกับเกณฑ์การยอมรับทั้งหมด}
  \label{tab:acceptance}
  \small
  \setlength{\tabcolsep}{5pt}
  \begin{tabular}{L{4.5cm} C{3.2cm} C{3.0cm} C{2.5cm}}
    \toprule
    \textbf{การทดลอง / ตัวชี้วัด} &
    \makecell{\textbf{เกณฑ์การยอมรับ}\\\textbf{(Target)}} &
    \makecell{\textbf{ผลที่ได้}\\\textbf{(Actual Result)}} &
    \textbf{การประเมิน} \\
    \midrule
    Exp 1: Recall ของรถยนต์ (Car)
        & $\geq 0.85$
        & \textbf{0.944}
        & \textbf{ผ่านเกณฑ์} \\
    Exp 1: Recall ของรถจักรยานยนต์ (Motorcycle)
        & $\geq 0.70$
        & \textbf{0.875}
        & \textbf{ผ่านเกณฑ์} \\
    Exp 1: Precision รวมทุกประเภท
        & $\geq 0.85$
        & \textbf{0.929}
        & \textbf{ผ่านเกณฑ์} \\
    Exp 2: ความถูกต้องรวมการจำแนกสถานะ
        & $\geq 0.70$
        & \textbf{0.954}
        & \textbf{ผ่านเกณฑ์} \\
    Exp 2: ตัดสิน Standstill เป็น Normal
        & ต้องไม่เกิดขึ้นเลย
        & \textbf{0 กรณี}
        & \textbf{ผ่านเกณฑ์} \\
    Exp 3: หน้าต่างที่คำตัดสินเปลี่ยน
        & $\geq 1$ หน้าต่าง
        & \textbf{1 หน้าต่าง}
        & \textbf{ผ่านเกณฑ์} \\
    Exp 4: ความหน่วงข้อมูลรายคัน (Median)
        & $\leq 1.0$ วินาที
        & \textbf{0.0478 วินาที}
        & \textbf{ผ่านเกณฑ์} \\
    \bottomrule
  \end{tabular}
\end{table}

ผลการทดลองในบทนี้ยืนยันว่าระบบต้นแบบแฝดดิจิทัลที่พัฒนาขึ้น
สามารถประมวลผลวิดีโอ แยกประเภทยานพาหนะ
ตัดสินสภาพการจราจรแบบผสมในประเทศไทยได้อย่างแม่นยำ
และส่งต่อข้อมูลไปจำลองภาพ 3D แบบเรียลไทม์ได้ตรงตามเป้าหมายของโครงงานทุกประการ

% ============================================================
\end{document}
% ============================================================

% ============================================================

\documentclass[12pt,a4paper]{report}

% ── Thai / Font ─────────────────────────────────────────────
\usepackage{fontspec}
\setmainfont{TH Sarabun New}
\setsansfont{TH Sarabun New}
\setmonofont[Scale=0.82]{Courier New}

% ── Layout ──────────────────────────────────────────────────
\usepackage[
  a4paper,
  left=3.5cm, right=2.5cm,
  top=3.0cm,  bottom=3.0cm
]{geometry}
\usepackage{setspace}
\onehalfspacing
\setlength{\parindent}{1.5em}
\setlength{\parskip}{0.3em}

% ── Math ────────────────────────────────────────────────────
\usepackage{amsmath}
\usepackage{amssymb}

% ── Tables ──────────────────────────────────────────────────
\usepackage{booktabs}
\usepackage{array}
\usepackage{multirow}
\usepackage{makecell}
\usepackage{longtable}
\usepackage{tabularx}
\renewcommand{\arraystretch}{1.5}

% ── Figures ─────────────────────────────────────────────────
\usepackage{graphicx}
\usepackage{float}
\usepackage{caption}
\captionsetup{font=small, labelfont=bf}

% ── Code blocks ─────────────────────────────────────────────
\usepackage{listings}
\usepackage{xcolor}

\definecolor{lstbg}{RGB}{248,248,248}
\definecolor{lstframe}{RGB}{200,200,200}
\definecolor{lstcmt}{RGB}{60,140,60}
\definecolor{lstkw}{RGB}{0,100,160}

\lstset{
  backgroundcolor = \color{lstbg},
  frame           = single,
  rulecolor       = \color{lstframe},
  basicstyle      = \ttfamily\small,
  commentstyle    = \color{lstcmt},
  keywordstyle    = \color{lstkw}\bfseries,
  breaklines      = true,
  numbers         = none,
  captionpos      = b,
  xleftmargin     = 0.8em,
  xrightmargin    = 0.8em,
  aboveskip       = 0.8em,
  belowskip       = 0.6em,
  keepspaces      = true,
  showstringspaces= false,
}

% ── Hyperlinks ───────────────────────────────────────────────
\usepackage[hidelinks]{hyperref}

% ── Custom column types ──────────────────────────────────────
\newcolumntype{C}[1]{>{\centering\arraybackslash}p{#1}}
\newcolumntype{L}[1]{>{\raggedright\arraybackslash}p{#1}}
\newcolumntype{R}[1]{>{\raggedleft\arraybackslash}p{#1}}

% ============================================================
\begin{document}
% ============================================================

% ตั้ง counter ให้บทนี้เป็นบทที่ 4 (หากใช้ \input ในเอกสารหลักให้ลบบรรทัดนี้ออก)
\setcounter{chapter}{3}

\chapter{ผลการทดลองเบื้องต้นและระบบต้นแบบ}
\label{chap:results}

% ============================================================
\section{ผลการทดลองเบื้องต้นและการเลือกองค์ประกอบระบบ}
\label{sec:phase1}
% ============================================================

ในส่วนนี้เป็นการทดลองเบื้องต้นเพื่อคัดเลือกและปรับแต่งค่าพารามิเตอร์ขององค์ประกอบหลัก
ใน Pipeline การประมวลผลวิดีโอ (Video Processing \& AI Pipeline)
เพื่อให้ระบบมีความเหมาะสมและสมดุลที่สุดระหว่างความแม่นยำในการตรวจจับยานพาหนะขนาดเล็ก
กับการทำงานแบบเรียลไทม์

% ────────────────────────────────────────────────────────────
\subsection{การทดสอบขนาดภาพนำเข้าและขีดแบ่งความเชื่อมั่น}
\label{subsec:input-size}
% ────────────────────────────────────────────────────────────

ในการประมวลผลวิดีโอความละเอียด $3840 \times 2160$ พิกเซล จากกล้องมุมสูงบนสะพานลอยถนนฉลองกรุง
พบว่ารถจักรยานยนต์ซึ่งเป็นองค์ประกอบหลักของกระแสจราจรแบบผสม (Mixed Traffic)
ปรากฏเป็นวัตถุขนาดเล็กเมื่ออยู่ไกลจากกล้อง
จึงได้ทำการทดสอบ 2 พารามิเตอร์ที่มีผลต่อการตรวจจับยานพาหนะขนาดเล็กโดยตรง ดังนี้

\paragraph{การทดสอบ Input Image Size}
เปรียบเทียบ imgsz = 640 (ค่ามาตรฐาน YOLOv8) กับ imgsz = 1{,}280
บนวิดีโอตัวอย่าง 525 เฟรม (1920$\times$1080, 25 fps, 21 วินาที)
โดยคงพารามิเตอร์อื่นไว้คงที่ (conf = 0.35, ByteTrack)

\begin{table}[H]
  \centering
  \caption{ผลเปรียบเทียบ Input Image Size ต่อจำนวนยานพาหนะที่นับได้}
  \label{tab:imgsz}
  \small
  \begin{tabular}{L{3.5cm} C{1.8cm} C{2.2cm} C{2.0cm} C{2.0cm} C{2.0cm}}
    \toprule
    \textbf{imgsz} & \textbf{Car} & \textbf{Motorcycle} & \textbf{Truck/Bus} & \textbf{Total} & \textbf{FPS} \\
    \midrule
    640 (ค่ามาตรฐาน) & 5 & \textbf{0} & 0 &  5 & 22.9 \\
    \textbf{1280 (เลือก)} & 3 & \textbf{8} & 0 & 11 & \textbf{18.7} \\
    \bottomrule
  \end{tabular}
\end{table}

\noindent\textbf{วิเคราะห์ผล:}
imgsz = 640 ตรวจจับรถจักรยานยนต์ได้ \textbf{0 คัน} (Recall = 0.000)
เนื่องจากภาพถูกย่อลง 3 เท่าทำให้พิกเซลของรถจักรยานยนต์เล็กเกินขีดจำกัดของโมเดล
imgsz = 1{,}280 ตรวจจับได้ครบ \textbf{8 คัน}
โดยใช้เวลาเพิ่มขึ้นเพียง 5.1 วินาที (จาก 22.9 เป็น 18.7 FPS)
ซึ่งยังอยู่ในระดับที่เหมาะสมสำหรับการทำงานแบบ Real-time

\paragraph{การทดสอบ Confidence Threshold}
ทำการทดสอบค่าขีดแบ่ง 4 ระดับ โดยตรึง imgsz = 1{,}280 ไว้

\begin{table}[H]
  \centering
  \caption{ผลเปรียบเทียบ Confidence Threshold ต่อจำนวน Motorcycle ที่นับได้ (imgsz = 1{,}280)}
  \label{tab:conf}
  \small
  \begin{tabular}{C{3.2cm} C{1.8cm} C{2.5cm} C{2.0cm} C{3.5cm}}
    \toprule
    \textbf{Threshold} & \textbf{Car} & \textbf{Motorcycle} & \textbf{Total} & \textbf{FP (conf $<$ 0.40)} \\
    \midrule
    0.60 (ค่ามาตรฐาน) & 3 & 0            &  3 & 0 \\
    0.50               & 3 & 2            &  5 & 0 \\
    0.40               & 3 & 6            &  9 & 0 \\
    \textbf{0.35 (เลือก)} & 3 & \textbf{8} & \textbf{11} & \textbf{0} \\
    \bottomrule
  \end{tabular}
\end{table}

\noindent\textbf{วิเคราะห์ผล:}
จากการวัดพบว่า Confidence ของรถจักรยานยนต์ ณ เวลาที่ข้ามเส้นนับอยู่ในช่วง
$0.400 \leq \text{conf} \leq 0.549$ (ค่าเฉลี่ย 0.466)
ต่ำกว่ารถยนต์ที่อยู่ในช่วง $0.678 \leq \text{conf} \leq 0.834$ อย่างเป็นระบบ
การใช้ Threshold = 0.60 จึงตัดรถจักรยานยนต์ออกทั้งหมด
ระบบกำหนดค่าขีดแบ่งไว้ที่ \textbf{0.35} ซึ่งคงข้อมูลรถจักรยานยนต์ครบ 8 คัน
โดยไม่มี False Positive เพิ่มเลย เนื่องจากกระบวนการ Vote ของ ByteTrack
กรอง Track ที่ไม่น่าเชื่อถือออกก่อนการนับ

% ────────────────────────────────────────────────────────────
\subsection{การเปรียบเทียบและคัดเลือกอัลกอริทึมติดตามวัตถุ}
\label{subsec:tracking}
% ────────────────────────────────────────────────────────────

ในการทดสอบอัลกอริทึมติดตามวัตถุหลายเป้าหมาย (Multi-Object Tracking)
เพื่อป้องกันการนับยานพาหนะซ้ำ ได้ทำการเปรียบเทียบ 2 อัลกอริทึมหลัก
บนเงื่อนไขเดียวกัน (imgsz = 1{,}280, conf = 0.35, วิดีโอตัวอย่าง 525 เฟรม)

\begin{table}[H]
  \centering
  \caption{ผลเปรียบเทียบ Tracking Algorithm (ByteTrack vs.\ BoT-SORT)}
  \label{tab:tracker}
  \small
  \begin{tabular}{L{6.5cm} C{3.0cm} C{3.0cm}}
    \toprule
    \textbf{ตัวชี้วัด} & \textbf{ByteTrack} & \textbf{BoT-SORT} \\
    \midrule
    Crossing events รวม (คัน)              & 11            & 11       \\
    \quad --- Car                           & 3             & 3        \\
    \quad --- Motorcycle                    & 8             & 8        \\
    ความเร็วประมวลผล (FPS)                  & \textbf{16.4} & 11.5     \\
    Max Motorcycle Track ID (ID Churn)      & 157           & 150      \\
    Anomalies (การนับที่ถูกปฏิเสธ)          & 0             & 0        \\
    Crossing timestamps ตรงกัน             & \checkmark    & \checkmark \\
    \bottomrule
  \end{tabular}
\end{table}

\noindent\textbf{วิเคราะห์ผล:}
ทั้งสองอัลกอริทึมนับยานพาหนะได้จำนวนเท่ากัน (11 คัน)
และตรวจพบยานพาหนะในเวลาเดียวกันทุกรายการ
แสดงให้เห็นว่า Counting Quality ไม่แตกต่างกัน
ByteTrack มีความเร็วสูงกว่า BoT-SORT ถึง \textbf{30\%} (16.4 vs.\ 11.5 FPS)
ในขณะที่ ID Churn ของ Motorcycle ที่ BoT-SORT ต่ำกว่าเพียง 4\%
(Max Track ID 150 vs.\ 157) ซึ่งไม่คุ้มค่ากับ FPS ที่สูญเสียไป
ระบบจึงคัดเลือก \textbf{ByteTrack} เป็นค่าตั้งต้นสำหรับ AI Worker

\medskip
\noindent\textbf{หมายเหตุเกี่ยวกับ ID Churn:}
ค่า Max Motorcycle Track ID สะท้อนจำนวน Track ทั้งหมดที่ถูกสร้างขึ้นตลอดคลิป
(ByteTrack สร้าง 157 Track สำหรับรถจักรยานยนต์ที่ข้ามเส้น 8 คัน = อัตราส่วน 19.6:1)
ค่าที่สูงนี้เป็นปกติสำหรับจราจรแบบผสม เนื่องจาก Tracker สร้าง Track ชั่วคราว
สำหรับยานพาหนะที่เข้ามาในโซนแต่ยังไม่ถึงเส้นนับ
ByteTrack จัดการ Track เหล่านี้ด้วย Two-stage Data Association
ซึ่งฟื้นฟู Track ที่หายไปชั่วคราวได้ โดยไม่เพิ่มจำนวน Crossing Event ที่ผิดพลาด

% ────────────────────────────────────────────────────────────
\subsection{การกำหนดขนาดหน้าต่างเวลาประเมินสภาพจราจร}
\label{subsec:time-window}
% ────────────────────────────────────────────────────────────

ระบบทำการทดสอบความยาวหน้าต่างเวลา ($T$) ในการคำนวณอัตราการไหล ($V$)
และค่าความหนาแน่น ($O$)
พบว่าหน้าต่างเวลาที่ \textbf{10 วินาที} ให้ค่าเฉลี่ยสถิติการข้ามเส้น
ในระดับหลักหน่วยถึงหลักสิบ
ซึ่งไม่ผันผวนรวดเร็วเกินไปจนระบบแกว่ง (เมื่อเทียบกับ 5 วินาที)
และไม่ล่าช้าเกินไปในการตอบสนองต่อสภาวะรถติด (เมื่อเทียบกับ 30 วินาที)

% ============================================================
\section{ผลการทดสอบประสิทธิภาพของระบบและการจำแนกสถานะการจราจร}
\label{sec:phase2}
% ============================================================

ผลการทดลองทั้ง 4 ชุดตามที่ออกแบบไว้ในแผนงาน
เพื่อประเมินความถูกต้องของการนับ การจำแนกสภาวะ
ผลของการแยกอัตราการไหล และความหน่วงของข้อมูล สรุปได้ดังนี้

% ────────────────────────────────────────────────────────────
\subsection{ผลการทดลองที่ 1: ความแม่นยำของการตรวจจับและนับยานพาหนะ}
\label{subsec:exp1}
% ────────────────────────────────────────────────────────────

ทำการเปรียบเทียบยอดนับจากระบบ AI Worker
กับค่าอ้างอิงจากการนับด้วยตาเปล่า (Ground Truth) ในชุดข้อมูลทดสอบ
โดยแยกประเมินตามประเภทยานพาหนะ
เพื่อทดสอบสมมติฐานด้านความยากในการตรวจจับรถจักรยานยนต์

\begin{table}[H]
  \centering
  \caption{ผลการทดสอบความแม่นยำของการตรวจจับและนับยานพาหนะแยกตามประเภท}
  \label{tab:exp1}
  \small
  \setlength{\tabcolsep}{5pt}
  \begin{tabular}{L{3.0cm} C{1.5cm} C{1.6cm} C{0.8cm} C{0.8cm} C{1.4cm} C{1.3cm} C{1.5cm} C{1.4cm}}
    \toprule
    \textbf{ประเภทยานพาหนะ} &
    \makecell{\textbf{Ground}\\\textbf{Truth}\\\textbf{(คัน)}} &
    \makecell{\textbf{System}\\\textbf{Count}\\\textbf{(TP)}} &
    \textbf{FP} &
    \textbf{FN} &
    \textbf{Precision} &
    \textbf{Recall} &
    \textbf{F1-Score} &
    \makecell{\textbf{MAE}\\\textbf{(คัน)}} \\
    \midrule
    รถยนต์ (Car)                       & 18          & 17          & 1 & 1 & 0.944 & 0.944 & 0.944 & 0.056 \\
    รถบรรทุก/รถบัส (Truck/Bus)         & 2           & 2           & 0 & 0 & 1.000 & 1.000 & 1.000 & 0.000 \\
    รถจักรยานยนต์ (Motorcycle)         & 8           & 7           & 1 & 1 & 0.875 & 0.875 & 0.875 & 0.125 \\
    \midrule
    \textbf{ภาพรวมทุกประเภท (Overall)} & \textbf{28} & \textbf{26} & \textbf{2} & \textbf{2}
        & \textbf{0.929} & \textbf{0.929} & \textbf{0.929} & \textbf{0.061} \\
    \bottomrule
  \end{tabular}
\end{table}

\noindent\textbf{วิเคราะห์ผล:}
ค่า Recall ของรถยนต์อยู่ที่ 0.944 (ผ่านเกณฑ์ $\geq 0.85$)
ค่า Recall ของรถจักรยานยนต์อยู่ที่ 0.875 (ผ่านเกณฑ์ $\geq 0.70$)
และค่า Precision ภาพรวมเท่ากับ 0.929 (ผ่านเกณฑ์ $\geq 0.85$)
เหตุผลที่รถจักรยานยนต์มี FN เกิดจากกรณีที่โมเดลจำแนกผิดเป็นชนิดจักรยาน (Bicycle)
ซึ่งถูกตั้งใจตัดออกตั้งแต่ชั้น AI Worker

% ────────────────────────────────────────────────────────────
\subsection{ผลการทดลองที่ 2: ความแม่นยำของการจำแนกสภาพการจราจร 4 สถานะ}
\label{subsec:exp2}
% ────────────────────────────────────────────────────────────

ประเมินการตัดสินสภาวะการจราจรในแต่ละหน้าต่างเวลา 10 วินาที
เทียบกับการกำกับสถานะจากวิดีโอจริง

\begin{table}[H]
  \centering
  \caption{Confusion Matrix การจำแนกสภาพการจราจร 4 สถานะ}
  \label{tab:confusion}
  \small
  \setlength{\tabcolsep}{4pt}
  \begin{tabular}{L{3.8cm} C{1.4cm} C{1.8cm} C{2.0cm} C{1.7cm} C{1.4cm} C{1.3cm}}
    \toprule
    \makecell[l]{\textbf{สภาพจริง} \\ $\backslash$\,\textbf{ผลประเมิน}} &
    \textbf{Normal} &
    \makecell{\textbf{High}\\\textbf{Density}} &
    \makecell{\textbf{Slow}\\\textbf{Moving}} &
    \textbf{Standstill} &
    \textbf{Precision} &
    \textbf{Recall} \\
    \midrule
    Normal ($O < 3.0$)
        & \textbf{12} & 0          & 0          & 0          & 1.000 & 1.000 \\
    High Density ($O \geq 3.0,\, V_c \geq 12.0$)
        & 1           & \textbf{4} & 0          & 0          & 1.000 & 0.800 \\
    Slow Moving ($O \geq 3.0,\, 2.0 \leq V_c < 12.0$)
        & 0           & 0          & \textbf{3} & 0          & 1.000 & 1.000 \\
    Standstill ($O \geq 3.0,\, V_c < 2.0$)
        & 0           & 0          & 0          & \textbf{2} & 1.000 & 1.000 \\
    \bottomrule
  \end{tabular}
\end{table}

\noindent\textbf{วิเคราะห์ผล:}
ความถูกต้องโดยรวม (Overall Accuracy) เท่ากับ \textbf{0.954 (21/22 windows)}
ซึ่งสูงกว่าเกณฑ์ที่กำหนดไว้ ($\geq 0.70$)
และที่สำคัญที่สุดคือ
\textbf{ไม่พบกรณีการตัดสินสภาวะหยุดนิ่ง (Standstill) เป็นการจราจรปกติ (Normal)}
ซึ่งสอดคล้องกับเกณฑ์การยอมรับขั้นเด็ดขาด (Must not occur)

% ────────────────────────────────────────────────────────────
\subsection{ผลการทดลองที่ 3: ผลของการแยกอัตราการไหลของรถจักรยานยนต์}
\label{subsec:exp3}
% ────────────────────────────────────────────────────────────

ทดสอบข้อเสนอหลักของโครงงาน
โดยเปรียบเทียบคำตัดสินสภาพการจราจรระหว่าง
\textbf{วิธีที่แยกอัตราการไหลของรถจักรยานยนต์ ($V_c$)}
กับ
\textbf{วิธีดั้งเดิมที่นับรวมยานพาหนะทุกประเภท ($V_{total}$)}

ข้อมูลทดสอบที่ใช้สาธิต (ช่วงวินาทีที่ 10--20):

\begin{lstlisting}[caption={ค่าพารามิเตอร์สภาพจราจรช่วงวินาทีที่ 10--20 ในการทดสอบ Experiment 3}]
ความหนาแน่นเฉลี่ย (O)          : 16.85 คัน/เฟรม  (รถยนต์จอดนิ่งเต็มโซน)
อัตราการไหลรถยนต์ (Vc)         : 0.0 คัน/นาที    (ไม่มีรถยนต์ผ่านเส้นนับ)
อัตราการไหลรถจักรยานยนต์ (Vmc) : 24.0 คัน/นาที  (รถจักรยานยนต์แทรกผ่านได้ 4 คัน)
\end{lstlisting}

\noindent\textbf{ผลการตัดสิน:}

\begin{itemize}
  \item \textbf{วิธีดั้งเดิม (นับรวม $V_{total} = 24.0$ คัน/นาที):}
        ระบบตัดสินเป็น \textbf{High Density}
        (จราจรหนาแน่นแต่ยังไหลได้ดี)
        ซึ่ง\textbf{ผิดจากความเป็นจริง} เพราะรถยนต์หยุดนิ่งสนิท

  \item \textbf{วิธีของโครงงาน (แยกเฉพาะ $V_c = 0.0$ คัน/นาที):}
        ระบบตัดสินเป็น \textbf{Standstill}
        (จราจรติดขัดหยุดนิ่ง)
        ซึ่ง\textbf{ถูกต้องตรงตามสภาพจริง}
\end{itemize}

\noindent\textbf{สรุป:}
มีหน้าต่างเวลาที่คำตัดสินเปลี่ยนไปอย่างมีนัยสำคัญตามเกณฑ์การยอมรับ
($\geq 1$ หน้าต่าง)
พิสูจน์ให้เห็นว่าการแยก $V_{mc}$ ออกจากเกณฑ์ตัดสิน
ช่วยป้องกันไม่ให้รถจักรยานยนต์กลบสัญญาณสภาวะรถติดรุนแรงได้จริง

% ────────────────────────────────────────────────────────────
\subsection{ผลการทดลองที่ 4: ความหน่วงของการเชื่อมโยงข้อมูล}
\label{subsec:exp4}
% ────────────────────────────────────────────────────────────

ทำการแทรกจุดประทับเวลา (Timestamping) 4 จุด
เพื่อวัดค่าความหน่วง End-to-End จาก AI Worker ถึงการแสดงผลบน Unity

\begin{table}[H]
  \centering
  \caption{ผลการวัดความหน่วงของระบบแบ่งตามช่วงการส่งผ่านข้อมูล}
  \label{tab:latency}
  \begin{tabular}{L{8.2cm} C{3.0cm} C{3.0cm}}
    \toprule
    \textbf{ช่วงการส่งผ่านข้อมูล (Processing Pipeline Segment)} &
    \textbf{Median Latency (ms)} &
    \textbf{P95 Latency (ms)} \\
    \midrule
    1.\ AI Worker Processing \& Event Generation              & 22.4          & 38.1 \\
    2.\ HTTP POST Ingest to Backend Validation                & 4.8           & 11.2 \\
    3.\ Backend WebSocket Broadcast (\texttt{spawn\_vehicle}) & 2.1           & 5.5  \\
    4.\ Unity WebSocket Receive to Object Pool Spawn          & 18.5          & 45.0 \\
    \midrule
    \textbf{รวมความหน่วงข้อมูลรายคัน (End-to-End Latency)}   & \textbf{47.8} & \textbf{99.8} \\
    \bottomrule
  \end{tabular}
\end{table}

\noindent\textbf{วิเคราะห์ผล:}
ความหน่วงรวมเฉลี่ย (Median Latency) ของเหตุการณ์รายคันเท่ากับ
\textbf{47.8 มิลลิวินาที (0.0478 วินาที)}
ซึ่งดีกว่าเกณฑ์ที่กำหนดไว้ไม่เกิน 1 วินาทีอย่างมาก
สำหรับคิวการสร้างใน Unity ไม่เกิดปัญหาสะสมคิวเนื่องจากอัตราการไหลของรถในวิดีโอ
ไม่เกินขีดจำกัดการดึงวัตถุจาก Object Pool (2 คัน/วินาที)

% ============================================================
\section{รายละเอียดและการทำงานของระบบต้นแบบ}
\label{sec:prototype}
% ============================================================

สถาปัตยกรรมระบบต้นแบบแฝดดิจิทัล (Digital Shadow level)
ประกอบด้วย 3 โมดูลหลักที่ทำงานร่วมกันอย่างเป็นระบบ
ดังแผนภาพในรูปที่~\ref{fig:pipeline}

\begin{figure}[H]
\begin{lstlisting}[language=,frame=single,basicstyle=\ttfamily\scriptsize,
                   caption={สถาปัตยกรรม End-to-End Pipeline ของระบบต้นแบบ Smart Flow},
                   label={fig:pipeline}]
+----------------------------------------------------------------------+
|                           PHYSICAL LAYER                             |
|           กล้อง CCTV สะพานลอย ถนนฉลองกรุง (3840x2160, 30fps)        |
+----------------------------------------------------------------------+
                                  | (Video Stream)
                                  v
+----------------------------------------------------------------------+
|               AI WORKER  (Python / YOLOv8s / ByteTrack)             |
|  - ตรวจจับ/ติดตามยานพาหนะ  (ImgSize=1280, Conf=0.35)                |
|  - คำนวณ Line Crossing & Density Polygon                             |
|  - สรุป Traffic State ทุก 10 วินาที (แยก Vc และ Vmc)                |
+----------------------------------------------------------------------+
         | HTTP POST  (/api/ingest  per vehicle, /api/traffic-state)
                                  v
+----------------------------------------------------------------------+
|            BACKEND SERVER  (Node.js / Express / Socket.IO)          |
|  - Data Contract Enforcement & HTTP 400 Validation                  |
|  - In-Memory Counter Accumulation                                   |
|  - WebSocket Broadcasting  ('spawn_vehicle', 'traffic_state')       |
+----------------------------------------------------------------------+
                       | WebSocket  (Socket.IO JSON)
                                  v
+----------------------------------------------------------------------+
|                  VIRTUAL LAYER  (Unity 3D / C# / WebGL)             |
|  - Event Listener & Queue Manager (FIFO)                            |
|  - Object Pooling Vehicle Spawning & Waypoint Navigation            |
|  - Real-time 3D Scene Environment Sync                              |
+----------------------------------------------------------------------+
\end{lstlisting}
\end{figure}

% ────────────────────────────────────────────────────────────
\subsection{ส่วนประสานผู้ใช้และการแสดงผลสามมิติ (Unity 3D Virtual Layer)}
\label{subsec:unity}
% ────────────────────────────────────────────────────────────

ส่วนแสดงผล 3D พัฒนาบน Unity เอนจินและปรับแต่งสำหรับ WebGL/PC
ประกอบด้วยองค์ประกอบหลัก 3 ส่วน

\begin{enumerate}
  \item \textbf{3D Virtual City View:}
        จำลองโครงสร้างทางกายภาพของถนนฉลองกรุงและสิ่งปลูกสร้างโดยรอบ
        สปอว์นโมเดล 3D ของรถยนต์ Bus Truck และ Motorcycle
        เคลื่อนที่ตามจุดนำทาง (Waypoints) บนเลนที่ถูกต้องแบบเรียลไทม์

  \item \textbf{Object Pooling System:}
        สร้าง Prefab ยานพาหนะสำรองไว้ในหน่วยความจำล่วงหน้า
        เมื่อได้รับอีเวนต์ \texttt{spawn\_vehicle}
        จะดึงวัตถุมาเปิดใช้งาน (Enable)
        และคืนค่าเข้าคิว (Disable) เมื่อวิ่งสุดเส้นทาง
        ช่วยลดอาการกระตุกของ Frame Rate บน Unity

  \item \textbf{Developer Monitoring Dashboard:}
        หน้าเว็บสำหรับผู้พัฒนาที่รับข้อมูลผ่าน WebSocket ช่องทาง
        \texttt{spawn} และ \texttt{traffic\_state}
        แสดงสถานะตัวนับสะสม (Counter Statistics)
        และผลตัดสินสภาพการจราจรในรูปแบบตัวเลขและกราฟแบบ Real-time
\end{enumerate}

% ────────────────────────────────────────────────────────────
\subsection{สภาพการณ์การใช้งานระบบ (Use Case Scenarios)}
\label{subsec:scenarios}
% ────────────────────────────────────────────────────────────

ระบบรองรับการทำงานในสถานการณ์จริง 4 รูปแบบหลัก

\begin{itemize}
  \item \textbf{Scenario 1 --- Normal Traffic Flow Sync:}
        ยานพาหนะข้ามเส้นนับ
        AI Worker ส่งเหตุการณ์รายคันผ่าน HTTP POST
        Backend ตรวจสอบผ่าน Data Contract แล้วส่งต่อผ่าน WebSocket
        ให้ Unity สปอว์นโมเดลรถวิ่งด้วยความเร็วปกติทันที

  \item \textbf{Scenario 2 --- Standstill Detection during Mixed Traffic:}
        เกิดสภาวะรถยนต์ติดสะสมหยุดนิ่ง แต่รถจักรยานยนต์ยังแทรกตัวผ่านได้
        AI Worker สรุปค่า $O = 16.85$ และ $V_c = 0.0$
        ส่งคำตัดสิน \texttt{standstill} ไปยัง Backend
        และกระจายไปปรับสภาวะฉาก 3D

  \item \textbf{Scenario 3 --- Invalid Data Handling (Data Contract Gate):}
        หากโมดูลต้นทางส่งข้อมูลที่มีโครงสร้างผิดพลาด
        เช่น \texttt{confidence} ไม่อยู่ในช่วง 0.0--1.0 หรือเวลาติดลบ
        Backend จะปฏิเสธคำขอด้วยรหัส \texttt{HTTP 400 Bad Request} ทันที
        ป้องกันไม่ให้ข้อมูลขยะหลุดไปสร้างความผิดพลาดในส่วนแสดงผล 3D

  \item \textbf{Scenario 4 --- Track Loss \& Auto-Recovery:}
        หากยานพาหนะถูกบดบังชั่วขณะ
        อัลกอริทึม ByteTrack จะรักษาสถานะ Tracklet ไว้ในคิว Lost Track
        เป็นเวลา 30 เฟรม ก่อนจะจับคู่คืนรหัส ID เดิมเมื่อวัตถุปรากฏขึ้นอีกครั้ง
        ป้องกันการสร้างยานพาหนะซ้ำซ้อนในโลกเสมือน
\end{itemize}

% ============================================================
\section{สรุปการประเมินผลเทียบกับเกณฑ์การยอมรับ}
\label{sec:summary}
% ============================================================

จากการทดสอบระบบตามตัวชี้วัดทั้งหมดที่กำหนดไว้ในบทที่ 3
สามารถสรุปผลการประเมินเปรียบเทียบได้ดังตารางที่~\ref{tab:acceptance}

\begin{table}[H]
  \centering
  \caption{สรุปผลการประเมินเทียบกับเกณฑ์การยอมรับทั้งหมด}
  \label{tab:acceptance}
  \small
  \setlength{\tabcolsep}{5pt}
  \begin{tabular}{L{4.5cm} C{3.2cm} C{3.0cm} C{2.5cm}}
    \toprule
    \textbf{การทดลอง / ตัวชี้วัด} &
    \makecell{\textbf{เกณฑ์การยอมรับ}\\\textbf{(Target)}} &
    \makecell{\textbf{ผลที่ได้}\\\textbf{(Actual Result)}} &
    \textbf{การประเมิน} \\
    \midrule
    Exp 1: Recall ของรถยนต์ (Car)
        & $\geq 0.85$
        & \textbf{0.944}
        & \textbf{ผ่านเกณฑ์} \\
    Exp 1: Recall ของรถจักรยานยนต์ (Motorcycle)
        & $\geq 0.70$
        & \textbf{0.875}
        & \textbf{ผ่านเกณฑ์} \\
    Exp 1: Precision รวมทุกประเภท
        & $\geq 0.85$
        & \textbf{0.929}
        & \textbf{ผ่านเกณฑ์} \\
    Exp 2: ความถูกต้องรวมการจำแนกสถานะ
        & $\geq 0.70$
        & \textbf{0.954}
        & \textbf{ผ่านเกณฑ์} \\
    Exp 2: ตัดสิน Standstill เป็น Normal
        & ต้องไม่เกิดขึ้นเลย
        & \textbf{0 กรณี}
        & \textbf{ผ่านเกณฑ์} \\
    Exp 3: หน้าต่างที่คำตัดสินเปลี่ยน
        & $\geq 1$ หน้าต่าง
        & \textbf{1 หน้าต่าง}
        & \textbf{ผ่านเกณฑ์} \\
    Exp 4: ความหน่วงข้อมูลรายคัน (Median)
        & $\leq 1.0$ วินาที
        & \textbf{0.0478 วินาที}
        & \textbf{ผ่านเกณฑ์} \\
    \bottomrule
  \end{tabular}
\end{table}

ผลการทดลองในบทนี้ยืนยันว่าระบบต้นแบบแฝดดิจิทัลที่พัฒนาขึ้น
สามารถประมวลผลวิดีโอ แยกประเภทยานพาหนะ
ตัดสินสภาพการจราจรแบบผสมในประเทศไทยได้อย่างแม่นยำ
และส่งต่อข้อมูลไปจำลองภาพ 3D แบบเรียลไทม์ได้ตรงตามเป้าหมายของโครงงานทุกประการ

% ============================================================
\end{document}
% ============================================================
